\documentclass[11pt,a4paper]{article}
% XeLaTeX-friendly preamble: use fontspec for Unicode fonts
\usepackage{fontspec}
\setmainfont{TeX Gyre Termes}
\usepackage[brazil]{babel}
\usepackage{geometry}
\geometry{margin=2.5cm}
\usepackage{amsmath,amssymb}
\usepackage{enumitem}
\usepackage{hyperref}
\usepackage{longtable}
\usepackage{graphicx}
\usepackage{float}
\usepackage{booktabs}
\usepackage{microtype}
\title{Revisão para a prova — Engenharia de Reservatórios I}
\author{Compilado}
\date{Gerado em 2026-04-22}
\begin{document}
\maketitle
\tableofcontents
\clearpage

% ---- file: cap1_transcription.md
\section{ENGENHARIA DE RESERVATÓRIOS I}
 \textbf{3º Ano – 2º Semestre — 2025–2026}

 \textbf{Docente:} Prof. Dr. Geraldo A. R. Ramos

 ---

\subsection{Sumário}
\begin{itemize}
\item 1. Sistema de Produção
\item 2. Engenharia de Reservatórios: funções e objetivos
\item 3. História da Engenharia de Reservatórios
\item 4. Tipos de reservatórios (diagramas de fases)
\item 5. Mecanismos de produção
\item 6. Consolidação da aula (exercícios)
\item 7. Tarefas
\item 8. Trabalho investigativo
\item 9. Bibliografia
\end{itemize}

 ---

\subsection{Objetivo geral}
 Analisar os princípios fundamentais da Engenharia de Reservatórios, abrangendo sistemas de produção, funções e objetivos da disciplina, evolução histórica, classificação de reservatórios segundo diagramas de fases e mecanismos de produção, visando a capacitação para integração de informações e apoio à tomada de decisão técnica em campos petrolíferos.

\subsubsection{Objetivos específicos}
\begin{itemize}
\item Descrever os componentes e a estrutura de um sistema de produção de hidrocarbonetos.
\item Explicar as principais funções e objetivos da Engenharia de Reservatórios na exploração e desenvolvimento de campos petrolíferos.
\item Relatar a evolução histórica da disciplina, destacando marcos tecnológicos e metodológicos.
\item Classificar os tipos de reservatórios com base em diagramas de fases (óleo, gás, condensado).
\item Comparar os principais mecanismos de produção: drive por pressão, expansão de gás, influxo de água e outros processos naturais.
\end{itemize}

\subsection{Sistema de Produção}
 O sistema de produção é o conjunto de equipamentos destinados à coleta, separação, tratamento, estocagem e escoamento dos fluidos produzidos e movimentados no campo de petróleo ou gás natural. A produção ocorre em três etapas principais:

\subsubsection{Etapas}
\begin{itemize}
\item Recuperação: fluxo de fluidos no reservatório até o fundo do poço.
\item Elevação: fluxo do fundo do poço até a superfície (topo do poço), que pode ser natural ou artificial.
\item Escoamento: fluxo do topo do poço até o separador, passando por reguladores de fluxo.
\end{itemize}

\subsubsection{Elementos do sistema petrolífero}
\begin{itemize}
\item **Armadilha (trapa):** arranjo estrutural que permite a acumulação de hidrocarbonetos (dobras, falhas, fraturas).
\item **Rocha geradora:** rocha rica em matéria orgânica que gera hidrocarbonetos.
\item **Rocha selante:** rocha com baixa permeabilidade que impede a migração de hidrocarbonetos.
\item **Migração:** processo que mobiliza os hidrocarbonetos da zona geradora para o reservatório.
\item **Sincronismo:** coincidência temporal dos processos geológicos necessários para formar um sistema petrolífero.
\item **Soterramento e catagênese:** processos térmicos que transformam querogênio em óleo, condensado ou gás.
\end{itemize}

 \emph{Nota:} Querogênio é a fração insolúvel da matéria orgânica; catagênese refere-se ao craqueamento térmico das moléculas maiores em moléculas menores sob aumento de temperatura e pressão.

\subsection{Engenharia de Reservatórios — Definição}
 A Engenharia de Reservatórios é o ramo da Engenharia do Petróleo que estuda o fluxo de fluidos (hidrocarbonetos) no interior das rochas porosas, visando maximizar o fator de recuperação. É uma atividade multidisciplinar que envolve física, química, geologia, engenharia, matemática e ciência da computação.

\subsection{Consolidação da aula — Exercícios}
 Responda as questões abaixo e justifique suas escolhas.

\begin{enumerate}
\item O óleo produzido em reservatórios leves caracteriza-se por:
\begin{itemize}
\item A) Alta viscosidade e baixa mobilidade
\item B) Baixa densidade e alta mobilidade
\item C) Formação de gás livre em excesso
\item D) Saturação irreductível de água igual a zero
\item E) Exclusivamente líquido sem variação de densidade
\end{itemize}
\end{enumerate}

\begin{enumerate}
\item O gás produzido em reservatórios influencia a recuperação de óleo porque:
\begin{itemize}
\item A) Reduz a pressão do reservatório e facilita o fluxo de óleo
\item B) Aumenta a densidade do óleo
\item C) Impede totalmente a produção de água
\item D) É irrelevante para a eficiência de produção
\item E) A saturação de óleo se torna irrelevante
\end{itemize}
\end{enumerate}

\begin{enumerate}
\item Reservatórios com mecanismo de water drive caracterizam-se por:
\begin{itemize}
\item A) Manutenção da pressão por influxo de água do aquífero
\item B) Elevação artificial de óleo via bombeio mecânico
\item C) Expansão de gás dissolvido no óleo
\item D) Produção exclusivamente de gás condensado
\item E) Inexistência de óleo saturado
\end{itemize}
\end{enumerate}

\begin{enumerate}
\item A viscosidade do óleo impacta diretamente:
\begin{itemize}
\item A) A taxa de fluxo pelo reservatório
\item B) A pressão do aquífero
\item C) A densidade da água
\item D) A compressibilidade da rocha
\item E) A saturação irreductível de gás
\end{itemize}
\end{enumerate}

\begin{enumerate}
\item Um mecanismo de produção por expansão de gás livre (gas-cap drive):
\begin{itemize}
\item A) Mantém pressão por expansão de gás acima do óleo
\item B) Depende exclusivamente de água injetada
\item C) Gera aumento da viscosidade do óleo
\item D) Reduz a saturação irreductível de água
\item E) Impede o fluxo de óleo
\end{itemize}
\end{enumerate}

\begin{enumerate}
\item O condensado produzido em reservatórios ocorre quando:
\begin{itemize}
\item A) O líquido condensado se separa do gás à medida que a pressão diminui
\item B) Óleo leve é produzido sem formação de gás
\item C) A água do reservatório se transforma em gás
\item D) A pressão está acima da pressão de saturação do gás
\item E) Não há variação de densidade com a pressão
\end{itemize}
\end{enumerate}

\begin{enumerate}
\item Qual das seguintes características NÃO define o mecanismo solution gas drive?
\begin{itemize}
\item A) Expansão do gás dissolvido no óleo
\item B) Redução gradual da pressão do reservatório
\item C) Formação de gás livre a partir do óleo
\item D) Dependência da pressão inicial para mobilidade
\item E) Produção contínua de óleo
\end{itemize}
\end{enumerate}

\begin{enumerate}
\item Em um reservatório com influxo de água (water influx), a produção de óleo é sustentada porque:
\begin{itemize}
\item A) A água do aquífero desloca o óleo em direção aos poços
\item B) A expansão do gás acima do óleo aumenta a viscosidade
\item C) Poços artificiais elevam a pressão
\item D) O óleo condensado impede a produção de água
\item E) A saturação irreductível de óleo é zero
\end{itemize}
\end{enumerate}

\begin{enumerate}
\item A produção de óleo em um reservatório undersaturated depende principalmente de:
\begin{itemize}
\item A) Pressão inicial do reservatório e mobilidade do óleo
\item B) Expansão do gás acima do óleo
\item C) Injeção de água ou gás
\item D) Formação de condensado
\item E) Elevação artificial exclusivamente
\end{itemize}
\end{enumerate}

\begin{enumerate}
\item O que caracteriza um reservatório com gas-cap drive?
\begin{itemize}
\item A) Gás livre no topo do reservatório que ajuda a manter a pressão
\item B) Água injetada artificialmente para suporte de pressão
\item C) Óleo pesado não fluido
\item D) Produção apenas de condensado líquido
\item E) Reservatório totalmente saturado sem gás
\end{itemize}
\end{enumerate}

\subsection{Tarefas}
\begin{enumerate}
\item Explique a diferença entre Sistema de Produção, Sistema Petrolífero e cadeia produtiva (cadeia de valor) da indústria petrolífera.
\item Faça uma ilustração no envelope de fases com os elementos mencionados no item 2.
\item Investigue as três principais teorias sobre a origem dos hidrocarbonetos:
\begin{itemize}
\item a) Qual a diferença entre as teorias?
\item b) Qual das teorias é mais aceita?
\item Qual é a área de atuação do Engenheiro de Reservatórios? Justifique.
\item Complete: HIDROCARBONETOS + NÃO HIDROCARBONETOS = _____________
\item Fundamente as respostas corretas e incorretas dos exercícios de consolidação.
\item Explique as principais propriedades dos fluidos produzidos em reservatórios e como influenciam a produção.
\item Descreva os mecanismos naturais de produção de um reservatório.
\item Compare os efeitos da expansão de gás e do influxo de água sobre a recuperação de óleo.
\item Discuta a importância de integrar informações sobre fluidos e mecanismos de produção no planejamento de campo.
\item Preencha a tabela sobre tipos de reservatórios, mecanismos de produção, condições iniciais, comportamento com o declínio de pressão e hidrocarbonetos produzidos.
\end{itemize}
\end{enumerate}

\subsection{Trabalho investigativo — História da Engenharia de Reservatórios}
 Os estudantes deverão elaborar um trabalho investigativo sobre a temática “História da Engenharia de Reservatórios”. O trabalho deve:
\begin{itemize}
\item Ser estruturado segundo as normas de um TCC: capa, resumo, sumário, introdução, desenvolvimento, conclusão e referências bibliográficas.
\item Ser produzido e apresentado em LaTeX, garantindo formatação acadêmica adequada.
\item Seguir rigorosamente as normas de citação e referência bibliográfica recomendadas pelo curso.
\end{itemize}

 \textbf{Data de submissão:} 24 de março de 2026.

\subsection{Definições importantes}
\begin{itemize}
\item **Sistema de Produção:** conjunto de instalações que permite a recuperação, elevação e escoamento dos fluidos produzidos desde o reservatório até instalações de superfície.
\item **Armadilha (trapa):** estrutura geológica que impede a migração de hidrocarbonetos e permite sua acumulação.
\item **Porosidade (\phi):** fração volumétrica de vazios na rocha, expressa como percentagem ou fração (por exemplo, 0.18 = 18%).
\item **Permeabilidade (k):** medida da facilidade com que um fluido atravessa a rocha (unidade: Darcy ou m²).
\item **Fator de volume de formação (B_o):** relação entre volume de óleo no reservatório e volume correspondente à superfície (Bo = V_reservoir_oil / V_surface_oil).
\item **Razão de solubilidade (R_s):** quantidade de gás dissolvido por unidade de óleo (usada em m3/m3 ou scf/STB).
\end{itemize}

\subsection{Exemplo prático — cálculo rápido de OOIP (fórmula de campo)}
Considere: Área = 100 acres; espessura neta h = 20 ft; porosidade \(\phi = 0.15\); saturação irreducível de água \(S_{wi}=0.25\); \(B_o = 1.2\) bbl/STB.

Usando a fórmula prática:
$$
OOIP = \frac{7758\,A\,h\,\phi\,(1-S_{wi})}{B_o}
$$
Substituindo valores:
$$
OOIP = \frac{7758 \times 100 \times 20 \times 0.15 \times (1-0.25)}{1.2} \approx 1{,}455\times10^{6}\;\text{STB}
$$
Interpretação: aproximadamente 1.45 milhões de barris em lugar antes da produção.

\subsection{Glossário de símbolos (rápido)}
\begin{itemize}
\item `A`: área (acres ou m²)
\item `h`: espessura (ft ou m)
\item `\phi`: porosidade (fração)
\item `S_w`, `S_{wi}`: saturação de água, saturação irreducível de água
\item `B_o`, `B_g`: fatores de volume de formação (óleo, gás)
\item `N`: óleo originalmente em lugar (OOIP)
\item `N_p`, `G_p`: produção acumulada de óleo e gás
\item `R_s`: razão gás/óleo dissolvido
\end{itemize}

\subsection{Dicas rápidas de estudo}
\begin{itemize}
\item Extraia os símbolos das equações e mantenha uma tabela pessoal de unidades (campo vs SI).
\item Para exercícios volumétricos, verifique unidades antes de aplicar fatores práticos (7758, 43560, etc.).
\item Use pequenos exemplos numéricos para validar o entendimento antes de resolver casos maiores.
\end{itemize}

---

Muito obrigada!


% ---- file: cap2_transcription.md


\section{ENGENHARIA DE RESERVATÓRIOS I}
\textbf{3º Ano – 2º Semestre — 2025–2026}

\textbf{Docente:} Prof. Dr. Geraldo A. R. Ramos

---

\subsection{Sumário}
\begin{itemize}
\item 2.1 Introdução
\item 2.2 Misturas e soluções
\item 2.3 Propriedades básicas dos fluidos
\item 2.3.1 Propriedades de óleo
\item 2.3.2 Propriedades dos gases
\item 2.3.3 Propriedades de misturas de hidrocarbonetos
\item 2.4 Consolidação da aula
\item 2.5 Tarefas
\item Bibliografia
\end{itemize}

---

\subsection{Objetivo geral}
Analisar as propriedades dos fluidos de reservatório por meio da caracterização de misturas e soluções, aplicação de equações de estado, interpretação de dados PVT e utilização de correlações empíricas, visando prever o comportamento dos fluidos e apoiar a tomada de decisão na engenharia de reservatórios.

\subsubsection{Objetivos específicos}
\begin{itemize}
\item Identificar os tipos de misturas e soluções presentes nos fluidos.
\item Descrever propriedades físicas e termodinâmicas relevantes (densidade, viscosidade, compressibilidade, fator de volume, etc.).
\item Aplicar equações de estado e correlações empíricas para estimativas de propriedades.
\item Calcular o fator volume de formação (B_o), a razão de solubilidade (R_s) e interpretar resultados PVT.
\item Avaliar a influência da viscosidade e do grau API no escoamento.
\end{itemize}

\subsection{2.1 Introdução}
Um fluido é uma substância que se deforma continuamente (escoa) sob ação de uma força tangencial. Para o engenheiro de reservatórios, conhecer as propriedades físico‑químicas e o comportamento de fase é essencial para estimar volumes in situ, planejar produção e definir estratégias de recuperação.

\subsubsection{Amostragem e análises}
\begin{itemize}
\item Amostragem em superfície (separadores e frascos de amostra).
\item Amostragem em fundo de poço (downhole sampling).
\item Amostragem em teste de formação (formation testing).
\end{itemize}

Análises normalmente realizadas:
\begin{itemize}
\item Determinação da composição (frações de hidrocarbonetos C1–Cn).
\item Análises PVT: ensaios CCE (constant composition expansion), DL (differential liberation), flash, etc.
\item Determinação de propriedades tecnológicas: $B_o$, $R_s$, viscosidade, densidade, Z‑factor.
\item Testes de compatibilidade de fluidos e análise de contaminantes.
\end{itemize}

\subsection{2.2 Misturas e soluções}
\begin{itemize}
\item Mistura: combinação física de componentes (gás + óleo + água) onde as fases podem coexistir.
\item Solução: situação em que um componente está dissolvido (ex.: gás dissolvido no óleo). A razão gás/óleo dissolvido é expressa por $R_s$ (m³ de gás por m³ de óleo, ou SCCM/Sm³ conforme unidade).
\item Pontos relevantes: ponto de bolha (bubble point) e ponto de orvalho (dew point) que definem o início da formação de uma fase livre.
\end{itemize}

\subsection{2.3 Propriedades básicas dos fluidos}

\subsubsection{2.3.1 Propriedades de óleo}
\begin{itemize}
\item Densidade e gravidade API: $API = \frac{141.5}{\rho_{sp}} - 131.5$ (forma prática para converter densidade específica em grau API).
\item Viscosidade: afeta diretamente a mobilidade e a taxa de fluxo através do meio poroso.
\item Fator volume de formação ($B_o$): relação entre volume no reservatório e volume à superfície; importante para conversão de volumes de reservatório para superficie.
\item Compressibilidade de óleo: influencia o declínio de pressão e estimativas volumétricas.
\end{itemize}

\subsubsection{2.3.2 Propriedades dos gases}
\begin{itemize}
\item Gases ideais obedecem à equação de estado $pV = nRT$ em condições de baixa pressão e alta temperatura.
\item Leis fundamentais: Boyle ( $V\propto 1/p$ ), Charles ( $V\propto T$ ), Avogadro.
\item Lei de Dalton (pressões parciais): a pressão total de uma mistura é a soma das pressões parciais.
\item Gases reais: desvio do comportamento ideal é quantificado pelo fator de compressibilidade $Z$ (tabelas/curvas de correção). Equações de estado cúbicas (Peng‑Robinson, Soave‑Redlich‑Kwong) são usadas em PVT.
\end{itemize}

\subsubsection{2.3.3 Propriedades de misturas de hidrocarbonetos}
\begin{itemize}
\item Comportamento de fases (monofásico vs bifásico) dependente de pressão e temperatura.
\item Razão gás‑óleo ($R_s$) e fator de formação $B_o$ em sistemas bifásicos.
\item Uso de análises PVT (flash/diferencial) para determinar curvas de $B_o(p)$, $R_s(p)$, viscosidade vs pressão e temperatura.
\end{itemize}

\subsection{2.4 Consolidação da aula}
Exemplos de perguntas para revisão:
\begin{itemize}
\item O que é $B_o$ e por que é importante para o cálculo de volumes produzidos?
\item Como o $R_s$ varia com a pressão e o que significa atingir o ponto de bolha?
\item Como a viscosidade do óleo influencia a mobilidade e o fator de recuperação?
\item Quando o gás pode ser tratado como ideal e quando deve‑se usar correções (Z‑factor)?
\end{itemize}

\subsection{2.5 Tarefas}
\begin{enumerate}
\item Desenhar o envelope de fases para uma mistura óleo‑gás e identificar ponto de bolha e ponto de orvalho.
\item Calcular $B_o$ e $R_s$ em um caso hipotético a partir de dados PVT simplificados (fornecerei os dados se desejar).
\item Fazer um resumo comparativo: gases ideais vs gases reais e apresentar quando usar equações de estado cúbicas.
\item Ler um ensaio PVT (CCE + DL) e resumir os principais resultados e curvas obtidas.
\end{enumerate}

\subsection{Definições importantes}
\begin{itemize}
\item **Fator de volume de formação ($B_o$):** razão entre o volume de óleo em condições de reservatório e o volume correspondente à superfície (bbl/STB ou m^3/m^3). Fórmula: $B_o = V_{res}/V_{surf}$.
\item **Razão de solubilidade ($R_s$):** quantidade de gás dissolvido por unidade de óleo (scf/STB ou m^3/m^3). No ponto de bolha, $R_s$ corresponde ao máximo gás dissolvido antes do aparecimento de gás livre.
\item **Fator de compressibilidade do gás ($Z$):** correção adimensional que quantifica o desvio do comportamento ideal do gás: $Z = pV/(nRT)$.
\item **Ponto de bolha / ponto de orvalho:** pressão (a uma dada temperatura) em que surge a primeira bolha de gás no líquido ou a primeira gota de líquido no gás.
\end{itemize}

\subsection{Exemplo prático — cálculo de $B_o$ e $R_s$}
Considere um ensaio PVT simplificado com resultados experimentais por 1 STB amostrado:
\begin{itemize}
\item Volume de óleo medido em condições de reservatório: $V_{res} = 1.20\;\text{bbl}$;
\item Volume correspondente à superfície: $V_{surf} = 1.00\;\text{STB}$;
\item Gás libertado no ensaio: $R_s = 400\;\text{scf/STB}$.
\end{itemize}

Cálculo de $B_o$:
$$
B_o = \frac{V_{res}}{V_{surf}} = \frac{1.20\;\text{bbl}}{1.00\;\text{STB}} = 1.20\;\text{bbl/STB}
$$
Interpretação: 1 STB de óleo de superfície corresponde a 1.20 bbl no reservatório.

Conversão de $R_s$ para unidades SI (opcional):
1 scf = 0.0283168 m^3; 1 STB = 0.1589873 m^3.
$$
R_s\,\text{(m}^3/\text{m}^3) = \frac{400\times0.0283168}{0.1589873} \approx 71.3\;\text{m}^3/\text{m}^3
$$

Observação: prefira as unidades pedidas no enunciado (scf/STB é muito comum). Converta unidades apenas quando necessário e mantenha consistência.

\subsection{Notas práticas e correlações}
\begin{itemize}
\item Quando não existirem dados experimentais, use correlações empíricas (p.ex. Standing, Vazquez‑Beggs) ou equações de estado cúbicas (Peng‑Robinson, SRK) para estimar $B_o$, $R_s$ e viscosidade.
\item Para gás, determine $Z(p,T)$ via charts ou equações para calcular volumes em condições padrão e avaliar expansibilidade.
\end{itemize}

\subsection{Glossário de símbolos (cap.2)}
\begin{itemize}
\item $B_o$: fator de volume de formação do óleo (bbl/STB ou m^3/m^3).
\item $R_s$: razão gás/óleo dissolvido (scf/STB ou m^3/m^3).
\item $Z$: fator de compressibilidade do gás (adimensional).
\item $\rho$: densidade (kg/m^3).
\item $\mu$: viscosidade (cP).
\end{itemize}

\subsection{Dicas rápidas de estudo}
\begin{itemize}
\item Verifique sempre as unidades antes de aplicar fórmulas; a maioria dos erros vem de conversões incorretas.
\item Plote curvas PVT simplificadas ($B_o(p)$, $R_s(p)$, $\mu(p)$) para compreender tendências com pressão.
\item Use exemplos numéricos pequenos (1 STB) para fixar interpretação física dos parâmetros.
\item Ao estudar equações de estado, foque primeiro no conceito físico (equilíbrio de fases) antes da implementação numérica.
\end{itemize}

---

Muito obrigado(a)!


% ---- file: cap3_transcription.md




of ⁨44⁩



Exercícios
\begin{enumerate}
\item Uma amostra saturada com óleo pesa 130 gramas. Sabe-se que a mesma amostra não saturada (seca)
pesa 105 gramas e a massa específica de óleo é de 840 kg/m3 =0,84g/cm3. O volume total de amostra
é de 180 cm3.
\item Determina a porosidade da amostra.\n\item Determina a porosidade a porosidade idealisada nas figuras seguintes:
c)
b)a)
ENGENHARIA DE RESERVATÓRIOS - I
GERALDO RAMOS, BSc, MSc, PhD
35
Exercícios
\item Estão disponíveis os seguintes dados:
V1 = 100cc
V2 = 100cc
p1 = 15,0 psi
p2 = 60,0 psi
\item Quando a válvula entre a Câmara 1 e a Câmara 2 é aberta, a pressão estabiliza\nem pf = 39,0 psi. Qual é o volume do grão da amostra do testemunho?
\item Determinar a porosidade média. Estão disponíveis os seguintes dados:
Amostra Porosidade, \%
1 10
2 12
3 11
4 13
5 14
6 10
7 17
ENGENHARIA DE RESERVATÓRIOS - I
GERALDO RAMOS, BSc, MSc, PhD
36
Exercícios
\item O peso seco de uma amostra é de 330 g, e seu peso quando saturada com água é de 360 g. O peso
aparente desta amostra em água é registado como 225 g. Dado que a densidade da água é 1 g / cm3,
determina a porosidade da amostra. Supondo que a amostra seja revestida com um material de peso
desprezível.
ENGENHARIA DE RESERVATÓRIOS - I
GERALDO RAMOS, BSc, MSc, PhD
\end{enumerate}
\section{ENGENHARIA DE RESERVATÓRIOS I}
\textbf{3º Ano – 2º Semestre — 2025–2026}

\textbf{Docente:} Prof. Dr. Geraldo A. R. Ramos

---

\subsection{Sumário}
\begin{itemize}
\item 3.1 Objectivo Geral e Objectivos Específicos
\item 3.2 Introdução
\item 3.3 Componentes da rocha
\item 3.4 Porosidade
\item 3.5 Compressibilidade
\item 3.6 Saturação de fluidos
\item 3.7 Permeabilidade
\item 3.8 Molhabilidade
\item 3.9 Capilaridade e pressão capilar
\item 3.10 Processos de embebição e drenagem
\item 3.11 Função J de Leverett
\item 3.12 Contacto hidrocarbonetos/água
\item 3.13 Consolidação da aula (Exercícios)
\item 3.14 Tarefas
\item Bibliografia
\end{itemize}

---

\subsection{3.1 Objectivo geral}
Descrever e caracterizar as propriedades das rochas de reservatório (porosidade, permeabilidade, compressibilidade, molhabilidade, capilaridade) e as relações entre esses parâmetros e o comportamento dos fluidos em meios porosos.

\subsubsection{Objectivos específicos}
\begin{itemize}
\item Identificar os componentes da rocha e os tipos de porosidade.
\item Medir e interpretar porosidade e permeabilidade de testemunhos.
\item Definir e aplicar conceitos de compressibilidade de rocha e poros.
\item Entender molhabilidade, capilaridade, embebição e drenagem.
\item Aplicar a função J de Leverett para correlacionar Pc entre reservatórios.
\end{itemize}

\subsection{3.2 Introdução}
A caracterização das rochas de reservatório é essencial para estimar volumes de hidrocarbonetos, modelar o fluxo de fluidos e projetar estratégias de produção e injecção. Ensaios laboratoriais em testemunhos (cores) fornecem propriedades fundamentais para o estudo do reservatório.

\subsection{3.3 Componentes da rocha}
\begin{itemize}
\item Matriz mineral (grãos) — determina a rigidez e a massa específica.
\item Poros — espaços que alojam fluidos; podem ser intergranulares, moldais, fraturais.
\item Cementos e matéria orgânica — alteram permeabilidade e porosidade.
\end{itemize}

\subsection{3.4 Porosidade}
\begin{itemize}
\item Porosidade total: fração volumétrica de vazios em relação ao volume total da rocha.
\item Porosidade efectiva: poros conectados que contribuem para o fluxo.
\item Métodos de medição: saturação/gravação de massas, porosimetria de mercúrio, tomografia, etc.
\end{itemize}

\subsection{3.5 Compressibilidade}
As rochas sob sobrecarga sofrem redução do volume poroso. A compressibilidade de poros (compressibilidade efectiva) é dada por:

$C_f = \dfrac{\Delta V_p / V_p}{\Delta p} = \dfrac{1}{V_p}\dfrac{\partial V_p}{\partial p}$

onde $C_f$ é a compressibilidade efectiva (unidade: psi⁻1 ou Pa⁻1), $V_p$ é o volume poroso e $\Delta p$ a variação de pressão.

Três compressibilidades relevantes:
\begin{itemize}
\item Compressibilidade da matriz: variação fraccional do volume sólido por variação de pressão.
\item Compressibilidade total da rocha: variação fraccional do volume total.
\item Compressibilidade dos poros (efectiva): variação fraccional do volume poroso.
\end{itemize}

\subsection{3.6 Saturação de fluidos}
\begin{itemize}
\item $S_w$: saturação de água; $S_{wc}$: saturação irreducível de água (connate). 
\item $S_o$: saturação de óleo; $S_{or}$: saturação residual de óleo.
\item Permeabilidades relativas dependem da saturação e do histórico de embebição/drenagem.
\end{itemize}

\subsection{3.7 Permeabilidade}
\begin{itemize}
\item Permeabilidade absoluta $k$ (unidade: Darcy ou m²) mede a facilidade de fluxo de um fluido.
\item Permeabilidade relativa $k_{r}$ descreve o fluxo de uma fase na presença de outras fases.
\end{itemize}

\subsection{3.8 Molhabilidade}
Molhabilidade é a tendência de um fluido aderir à superfície sólida na presença de outros fluidos; é caracterizada pelo ângulo de contacto $\theta$:

$\cos\theta = \dfrac{\sigma_{so}-\sigma_{sw}}{\sigma_{wo}}$

onde $\sigma_{so},\sigma_{sw},\sigma_{wo}$ são as tensões superficiais entre sólido‑óleo, sólido‑água e água‑óleo, respectivamente. A fase molhante tende a ocupar poros menores.

\subsection{3.9 Capilaridade e pressão capilar}
Pressão capilar: $P_c = P_{non-wetting} - P_{wetting}$. A curva $P_c(S_w)$ relaciona pressão capilar e saturação; é controlada por radio de poros effective, tensão interfacial e molhabilidade.

\subsection{3.10 Embebição e drenagem}
\begin{itemize}
\item Embebição: processo onde a fase molhante invade a rocha, reduzindo saturação da fase não‑molhante.
\item Drenagem: processo oposto, onde a fase não‑molhante invade a rocha.
\item Histerese: curvas de Pc e permeabilidade relativa dependem do caminho (embebição vs drenagem).
\end{itemize}

\subsection{3.11 Função J de Leverett}
A função de Leverett permite correlacionar curvas de pressão capilar entre rochas com propriedades diferentes:

$J(S_w)=\dfrac{P_c(S_w)\sqrt{\dfrac{k}{\phi}}}{\sigma\cos\theta}$

onde $k$ é a permeabilidade absoluta, $\phi$ a porosidade, $\sigma$ a tensão interfacial e $\theta$ o ângulo de contacto.

\subsection{3.12 Contacto hidrocarbonetos/água}
Em condições estáticas, as pressões nas interfaces obedecem ao equilíbrio hidrostático. Para o nível de óleo livre (FOL):

$p_{Z_{FOL}} = p_g + \rho_g g (Z_{FOL}-Z_g)$

e pelo lado do óleo:

$p_{Z_{FOL}} = p_o - \rho_o g (Z_o - Z_{FOL})$

Combinando as duas expressões obtém‑se a posição do nível em função das pressões e densidades (ver equações no material original).

\subsection{3.13 Consolidação da aula — Exercícios}
\begin{enumerate}
\item Uma amostra saturada com óleo pesa 130 g; seca pesa 105 g. Massa específica do óleo = 0,84 g/cm³. Volume total da amostra = 180 cm³. Determine a porosidade da amostra.
\end{enumerate}

\begin{enumerate}
\item Determine a porosidade idealizada nas figuras (a), (b) e (c). (Figuras não incluídas — desenhar/esboçar.)
\end{enumerate}

\begin{enumerate}
\item Dados: $V_1=100\,$cc; $V_2=100\,$cc; $p_1=15{\,}\text{psi}$; $p_2=60{\,}\text{psi}$. Ao abrir a válvula, a pressão final é $p_f=39{\,}\text{psi}$. Qual é o volume do grão da amostra do testemunho?
\end{enumerate}

\begin{enumerate}
\item Calcule a porosidade média para as amostras com porosidades: 10, 12, 11, 13, 14, 10, 17%.
\end{enumerate}

\begin{enumerate}
\item Peso seco = 330 g; peso saturado com água = 360 g; peso aparente em água = 225 g. Densidade da água = 1 g/cm³. Determine a porosidade (assuma revestimento desprezível).
\end{enumerate}

\begin{enumerate}
\item Complete a tabela de análises de rotina e testes complementares (porosidade, permeabilidade, saturações, litologia, permeabilidade vertical, core‑gamma, massa específica da matriz, propriedades físico‑químicas de óleo/água).
\end{enumerate}

\begin{enumerate}
\item Análise especial de testemunho — identifique testes estáticos (compressibilidade, estudo petrográfico, molhabilidade, capilaridade, acústica, elétrica) e dinâmicos (estudo do fluxo, testes de fluxo).
\end{enumerate}

\begin{enumerate}
\item Medição laboratorial da porosidade: descrever procedimentos para amostras medidas após 1, 2 e 3 semanas (hidratação, secagem, estabilização de massa, etc.).
\end{enumerate}

\subsection{3.14 Tarefas}
\begin{enumerate}
\item Explique o conceito de rocha de reservatório e sua importância.
\item Diferencie porosidade total e efectiva.
\item Analise os efeitos da compactação na porosidade.
\item Descreva a influência dos processos diagenéticos.
\item Fundamente as questões corretas e incorrectas dos exercícios de consolidação.
\end{enumerate}

\subsection{Nota sobre unidades — PSI}
PSI: pound‑force per square inch. Duas escalas comuns:

$\text{psia} = \text{psig} + 14.696$

onde 14.696 psi é aproximadamente a pressão atmosférica ao nível do mar.

\subsection{Definições importantes}
\begin{itemize}
\item **Porosidade ($\phi$):** fração volumétrica de vazios do volume total da rocha: $\phi = V_p / V_t$.
\item **Porosidade efectiva ($\phi_e$):** porosidade conectada que contribui para o fluxo.
\item **Permeabilidade ($k$):** facilidade de fluxo através da rocha; unidade prática: Darcy (D). É uma propriedade intrínseca da rocha.
\item **Volume poroso ($V_p$):** volume ocupado por fluidos nos poros; $V_p=\phi\,V_t$.
\item **Compressibilidade dos poros ($C_f$):** variação fraccional do volume poroso por variação de pressão: $C_f=\dfrac{1}{V_p}\dfrac{\partial V_p}{\partial p}$.
\end{itemize}

\subsection{Exemplos práticos — resolução (passo a passo)}

\subsubsection{Exemplo 1 — (Exercício 1)}
Dados: amostra saturada com óleo pesa $130\,$g; amostra seca pesa $105\,$g; massa específica do óleo $\rho_o=0{,}84\,$g/cm$^3$; volume total da amostra $V_t=180\,$cm$^3$.

\begin{enumerate}
\item Volume de fluido nos poros:
$$
V_f = \dfrac{m_{sat}-m_{dry}}{\rho_f} = \dfrac{130-105}{0{,}84} \approx 29{,}76\;\text{cm}^3.
$$
\item Porosidade:
$$
\phi = \dfrac{V_p}{V_t} = \dfrac{29{,}76}{180} \approx 0{,}165 \approx 16{,}5\%.
$$
\end{enumerate}

Conclusão: a porosidade da amostra é aproximadamente $16{,}5\%$.

\subsubsection{Exemplo 2 — (Exercício 5, método de Arquímedes)}
Dados: peso seco $m_{dry}=330\,$g; peso saturado $m_{sat}=360\,$g; peso aparente em água $m_{ap\_agua}=225\,$g; densidade da água $\rho_{agua}=1\,$g/cm$^3$.

\begin{enumerate}
\item Volume total da amostra via empuxo (diferença entre peso em ar e peso aparente em água):
$$
V_t = \dfrac{m_{sat}-m_{ap\_agua}}{\rho_{agua}} = \dfrac{360-225}{1} = 135\;\text{cm}^3.
$$
\item Volume poroso (água/fluido preenchendo os poros):
$$
V_p = m_{sat}-m_{dry} = 360-330 = 30\;\text{cm}^3.
$$
\item Porosidade:
$$
\phi = \dfrac{V_p}{V_t} = \dfrac{30}{135} \approx 0{,}222 = 22{,}2\%.
$$
\end{enumerate}

Interpretação: o método de Arquímedes é útil quando o volume total não é fornecido diretamente.

\subsection{Exemplo rápido — compressibilidade de poros}
Se um volume poroso $V_p=100\,$cm$^3$ diminui $0{,}2\,$cm$^3$ quando a pressão aumenta $\Delta p=100\,$psi, então a compressibilidade efectiva vale:
$$
C_f = \dfrac{\Delta V_p/V_p}{\Delta p} = \dfrac{0{,}2/100}{100} = 2\times10^{-5}\;\text{psi}^{-1}
$$
(apresente o valor em módulo como magnitude positiva quando reportar resultados).

\subsection{Medidas e técnicas comuns}
\begin{itemize}
\item Análise de testemunhos (core): porosidade por pesagem, porosimetria de mercúrio, perméabilímetro de bancada.
\item Técnicas não destrutivas: NMR, micro‑CT, imageamento 3D para distribuição de poros.
\item Medição de pressão capilar: centrífuga e porosimetria; integrar com curvas de permeabilidade relativa.
\item Integração: comparar core, registos de poço (neutron/densidade) e NMR para validação.
\end{itemize}

\subsection{Glossário de símbolos (cap.3)}
\begin{itemize}
\item $\phi$ — porosidade total (fração).
\item $\phi_e$ — porosidade efectiva (fração).
\item $V_p$ — volume poroso (cm$^3$ ou m$^3$).
\item $V_t$ — volume total da amostra/elemento (cm$^3$ ou m$^3$).
\item $S_w, S_o$ — saturações de água e óleo (frações adimensionais).
\item $k$ — permeabilidade absoluta (Darcy, m$^2$).
\item $k_r$ — permeabilidade relativa (adimensional).
\item $P_c$ — pressão capilar (Pa ou psi).
\item $C_f$ — compressibilidade efectiva dos poros (Pa$^{-1}$ ou psi$^{-1}$).
\end{itemize}

\subsection{Dicas rápidas de estudo}
\begin{itemize}
\item Anote sempre as unidades e converta antes de aplicar fórmulas (campo vs SI).
\item Modele pequenos exemplos numéricos (1 amostra; 1 STB) para fixar conceitos.
\item Compare resultados laboratoriais com registos de poço; investigue discrepâncias.
\item Para curvas $P_c(S_w)$, utilize a função de Leverett $J(S_w)$ para correlacionar entre rochas.
\end{itemize}

---

Muito obrigado(a)!

% ---- file: cap4_transcription.md



1



ENGENHARIA DE RESERVATÓRIOS I
3º ANO – 2º SEMESTRE
2025 - 2026
1
Docente: Prof. Dr. Geraldo A. R. Ramos
1Engenharia de Reservatórios I
Geraldo Ramos, BSc, MSc, PhD
Capítulo 4 - Integração de dados e Cálculo Volumétrico de Hidrocarbonetos
\section{ENGENHARIA DE RESERVATÓRIOS I}
 \textbf{3º Ano – 2º Semestre — 2025–2026}

 \textbf{Docente:} Prof. Dr. Geraldo A. R. Ramos

 ---

\subsection{Capítulo 4 — Integração de dados e Cálculo Volumétrico de Hidrocarbonetos}

\subsubsection{Sumário}
\begin{itemize}
\item 4.1 Objectivo Geral e Objectivos Específicos
\item 4.2 Prefixos
\item 4.3 Introdução
\item 4.4 Definições
\item 4.5 Métodos de cálculo de volumes originalmente existentes
\item 4.6 Método volumétrico (parâmetros e fórmulas)
\item 4.7 Integração de dados e incertezas
\item 4.8 Consolidação da aula (exercícios)
\item 4.9 Tarefas
\item Bibliografia
\end{itemize}

 ---

\subsection{4.1 Objetivo geral}
Aplicar integração de dados geológicos, petrofísicos e de engenharia para estimar volumes originais de hidrocarbonetos (OOIP / OGIP) e volumes recuperáveis, considerando incertezas e apresentando resultados úteis para tomada de decisão.

\subsubsection{Objetivos específicos}
\begin{itemize}
\item Definir conceitos fundamentais relacionados a volumes em reservatórios.
\item Explicar métodos de cálculo (volumétrico, material balance, decline, probabilístico).
\item Descrever parâmetros essenciais do método volumétrico e suas fontes de incerteza.
\item Demonstrar integração de dados (mapas de espessura, Net‑to‑Gross, phi, Sw, Bo, Bg).
\end{itemize}

\subsection{4.2 Prefixos e unidades (resumo)}
| Prefixo | Símbolo | Fator |
|---|---:|---:|
| Kilo | k | 10^3 |
| Mega | M | 10^6 |
| Giga | G | 10^9 |
| Milli | m | 10^-3 |
| Micro | µ | 10^-6 |

> Nota: escolha unidades consistentes (SI ou campo — acres/ft/bbl). Ao usar fórmulas práticas, aplique os fatores de conversão corretos.

\subsection{4.3 Introdução}
O método volumétrico estima o volume de fluido originalmente existente no reservatório a partir de propriedades petrofísicas e geométricas: área, espessura neta, porosidade, saturação, e fatores de volume de formação. É a primeira estimativa após descoberta e serve de base para cálculos de reservas.

\subsection{4.4 Definições essenciais}
\begin{itemize}
\item $V_r$: volume de rocha do reservatório (m³ ou ft³);
\item $\phi$: porosidade (fração);
\item $S_w$: saturação de água (fração); $S_o=1-S_w$ saturação de óleo;
\item $B_o, B_g$: fatores de volume de formação (reservatório → superfície);
\item $N$: número de barris padrão (STB) ou volume de gás em condições padrão.
\end{itemize}

\subsection{4.5 Métodos de cálculo de volumes originalmente existentes}
\begin{itemize}
\item Método volumétrico (petrofísico/geométrico).
\item Material balance (balanço de material) — requer dados de produção e pressão.
\item Métodos empíricos e baseados em declínio — para reservas quando há produção.
\item Abordagens probabilísticas (Monte Carlo) — para quantificar incerteza.
\end{itemize}

\subsection{4.6 Método volumétrico}
\subsubsection{4.6.1 Forma geral (reservoir conditions)}
O volume de hidrocarboneto no reservatório (volume de fluido em reservatório):
$$
V_{fluid,res} = V_r\,\phi\,(1-S_w)
$$

Para converter para unidades de superfície usa‑se o fator de volume de formação $B$:
$$
N_{surface} = \frac{V_{fluid,res}}{B}
$$

\subsubsection{4.6.2 Fórmulas práticas (exemplos)}
\begin{itemize}
\item OOIP (stock tank barrels, usando unidades campo):
$$
OOIP = \frac{7758\,A\,h\,\phi\,(1-S_w)}{B_o}
$$
onde $A$ em acres, $h$ em ft, $\phi$ fração, $B_o$ em bbl/STB.
\end{itemize}

\begin{itemize}
\item OGIP (scf ou Sm^3 dependendo da convenção):
$$
OGIP = \frac{43560\,A\,h\,\phi\,(1-S_w)}{B_g}
$$
\end{itemize}

> Observação: ao usar SI, converta $A$ (m²), $h$ (m) e fatores de conversão adequados.

\subsubsection{4.6.3 Estimativa de volume recuperável (reservas)}
Reservas (recuperáveis) podem ser expressas como função das condições iniciais e finais:

Para gás (exemplo conceitual):
$$
N_R = V_r\,\phi\,(1-S_{w_i})\left(\frac{1}{B_{g_i}} - \frac{1}{B_{g_a}}\right)
$$

Para óleo (contando saturações iniciais e residuais):
$$
N_R = V_r\,\phi\,(1-S_{w_i})\left(\frac{S_{o_i}}{B_{o_i}} - \frac{S_{o_r}}{B_{o_r}}\right)
$$

Onde índices $i$ e $a$ referem‑se às condições iniciais e de abandono; $S_{o_r}$ e $B_{o_r}$ são saturação de óleo residual e fator de volume no estado residual.

\subsubsection{4.6.4 Reserva por fator de recuperação}
Reservas também podem ser estimadas por:
$$
N_R = N \times F_R
$$
onde $N$ é o volume originalmente em lugar e $F_R$ o fator de recuperação (fração).

\subsection{4.7 Integração de dados e incertezas}
Passos práticos:
\begin{itemize}
\item QA/QC das fontes: mapas de batimetria, logs, cores, testes de poço, análises PVT.
\item Construir mapas de Net‑to‑Gross, porosidade média, Sw média por setor.
\item Calcular sensibilidade dos parâmetros-chave (A, h, φ, Sw, Bo) — análise P10/P50/P90.
\item Usar Monte Carlo para propagar incertezas e gerar distribuições de OOIP/OGIP e reservas.
\end{itemize}

Principais fontes de incerteza: interpretação de limites (cut‑offs), heterogeneidade vertical e lateral, erros de core/log, variabilidade PVT.

\subsection{4.8 Consolidação da aula — Exercícios}
\begin{enumerate}
\item Dado: Área = 200 acres; h_net = 30 ft; φ = 0.18; Swi = 0.25; Bo = 1.2 bbl/STB. Calcule OOIP (use fórmula prática).
\item Faça uma análise de sensibilidade variando φ entre 0.15–0.22 e interprete o impacto no OOIP.
\item Explique como Net‑to‑Gross e cut‑offs alteram a estimativa volumétrica.
\end{enumerate}

\subsection{4.9 Tarefas}
\begin{enumerate}
\item Reunir os dados petrofísicos e montar uma planilha de cálculo volumétrico para um caso hipotético.
\item Rodar uma simulação Monte Carlo simples (50–100 iterações) variando φ, Sw e Bo e apresentar P10/P50/P90.
\item Elaborar um pequeno relatório discutindo fontes de incerteza e medidas para redução de risco.
\end{enumerate}

\subsection{Definições importantes}
\begin{itemize}
\item **Net‑to‑Gross (N/G):** fração da coluna que é considerada 'net pay' (contribui para produção).
\item **Cut‑off:** valor mínimo de porosidade/permeabilidade/saturação usado para definir zonas produtivas.
\item **Fator de recuperação (F_R):** fração do volume originalmente em lugar que se espera recuperar.
\item **Área efetiva (A_eff):** área do reservatório considerada produtiva após aplicar cut‑offs e mapas de net.
\end{itemize}

\subsection{Exemplo prático — cálculo de OOIP (exercício 4.8)}
Dados: Área $A=200\,$acres; $h_{net}=30\,$ft; $\phi=0.18$; $S_{wi}=0.25$; $B_o=1.2\,$bbl/STB.

Usando a fórmula prática:
$$
OOIP = \frac{7758\,A\,h\,\phi\,(1-S_w)}{B_o}
$$
Substituindo os valores:
$$
OOIP = \frac{7758 \times 200 \times 30 \times 0.18 \times (1-0.25)}{1.2}
$$
Calculando passo a passo:
$$
7758\times200 = 1\,551\,600\\
1\,551\,600\times30 = 46\,548\,000\\
46\,548\,000\times0.18 = 8\,378\,640\\
8\,378\,640\times0.75 = 6\,283\,980\\
	ext{OOIP} = \dfrac{6\,283\,980}{1.2} \approx 5\,236\,650\;\text{STB}
$$
Interpretação: aproximadamente $5.24\times10^6$ STB (5,24 milhões de barris originalmente em lugar).

\subsection{Observações sobre unidades}
\begin{itemize}
\item Fórmulas práticas usam unidades de campo (acres, ft, bbl). Para SI, converta área e espessura e aplique os fatores apropriados.
\item Sempre verifique se $B_o$ e $B_g$ estão no mesmo sistema de unidades antes de dividir.
\end{itemize}

\subsection{Glossário de símbolos (cap.4)}
\begin{itemize}
\item `A` — área (acres ou m²).
\item `h` — espessura neta (ft ou m).
\item `\phi` — porosidade (fração).
\item `S_w` — saturação de água (fração); `S_o = 1-S_w`.
\item `B_o`, `B_g` — fatores de volume de formação do óleo e do gás.
\item `N`, `OGIP` — volumes originalmente em lugar (STB para óleo; scf ou Sm³ para gás).
\item `F_R` — fator de recuperação (fração).
\end{itemize}

\subsection{Dicas rápidas de estudo}
\begin{itemize}
\item Monte uma planilha com a fórmula prática e permita variação de parâmetros (sensibilidade).
\item Construa mapas de A_eff e h_net para setores diferentes e compare OOIP setorial.
\item Use Monte Carlo para obter P10/P50/P90 e comunicar risco na estimativa.
\end{itemize}

---

Muito obrigado(a)!

% ---- file: cap5_transcription.md



1



ENGENHARIA DE RESERVATÓRIOS I
3º ANO – 2º SEMESTRE
2025 - 2026
1
Docente: Prof. Dr. Geraldo A. R. Ramos
1Engenharia de Reservatórios I
Geraldo Ramos, BSc, MSc, PhD
Capítulo 5 – Equação de Balanço de Materiais
▪ 5.1. Objectivo Geral e Objectivos Específicos
▪ 5.2. Introdução à Equação de Balanço de Materiais
▪ 5.3. Equação de balanço de materiais Generalizada;
▪ 5.4. Equação de Balanço de Materiais em reservatórios de óleo;
▪ 5.5. Equação de Balanço de Materiais em reservatórios de gás;
▪ 5.6. Técnicas de análise das Equações de Balanço de Materiais
▪ 5.7. Vantagens e desvantagens das equações de balanço de materiais;
▪ .5.8. Consolidação da Aula
▪ 5.9. Tarefas
\section{ENGENHARIA DE RESERVATÓRIOS I}
 \textbf{3º Ano – 2º Semestre — 2025–2026}

 \textbf{Docente:} Prof. Dr. Geraldo A. R. Ramos

 ---

\subsection{Capítulo 5 — Equação de Balanço de Materiais}

\subsubsection{Sumário}
\begin{itemize}
\item 5.1 Objetivo geral e objetivos específicos
\item 5.2 Introdução à EBM
\item 5.3 Equação de Balanço de Materiais (forma geral)
\item 5.4 EBM para reservatórios de óleo (linearização)
\item 5.5 EBM para reservatórios de gás
\item 5.6 Técnicas de análise (gráficos e métodos)
\item 5.7 Vantagens e limitações
\item 5.8 Consolidação da aula (exercícios)
\item 5.9 Tarefas
\item Bibliografia
\end{itemize}

---

\subsection{5.1 Objetivo geral}
Aplicar a equação de balanço de materiais (EBM) para quantificar volumes em reservatórios de óleo e gás, integrando dados de produção, PVT e propriedades do reservatório, e avaliando a incerteza das estimativas.

\subsubsection{Objetivos específicos}
\begin{itemize}
\item Explicar os fundamentos e hipóteses da EBM.
\item Interpretar os termos da EBM generalizada.
\item Aplicar a linearização (Havlena & Odeh) para estimar volumes e parâmetros desconhecidos.
\item Comparar técnicas de análise e discutir vantagens/limitações.
\end{itemize}

\subsection{5.2 Introdução}
A EBM baseia‑se no balanço de massa dos fluidos presentes nos poros do reservatório. Em termos práticos, relaciona o fluido produzido e injetado com a variação do estoque no reservatório, permitindo estimar o volume originalmente em lugar (N) ou parâmetros como o fator de recuperação.

A técnica situa‑se entre a estimativa volumétrica e a simulação: usa dados de produção (e PVT) para inferir o estoque sem exigir dimensões detalhadas do reservatório.

\subsection{5.3 Equação de Balanço de Materiais — forma geral}
Em termos gerais, a EBM pode ser escrita como uma relação entre os termos de produção/injeção/influxo e as variações volumétricas e compressibilidades. Uma forma útil para análise prática é escrever a equação como:

$$
F = N E_o + m E_g + (1+m) E_{f,w} + W_e
$$

onde os termos são definidos abaixo (Havlena \& Odeh, 1963) e $F$ representa o lado conhecido (produções e injecções).

\subsection{5.4 EBM para reservatórios de óleo — linearização (Havlena & Odeh)}
Para facilitar a estimação dos parâmetros desconhecidos, a EBM é linearizada agrupando termos de variação de volume, compressibilidade e contribuição do gas‑cap. As definições práticas são:

$$E_o = B_o - B_{o,i} + (R_{s,i} - R_s)\,B_g$$

$$E_g = B_{o,i}\left(\frac{B_g}{B_{g,i}} - 1\right)$$

$$E_{f,w} = B_{o,i}\left(S_{w,i}c_w + \frac{c_f}{1-S_{w,i}}\Delta p\right)$$

e o termo conhecido (lado esquerdo) é, por exemplo:

$$F = N_p B_o + G_p - R_s B_g + W_p B_w - W_{inj} B_w - G_{inj} B_{g,inj}$$

Assim, a equação linear para regressão fica:

$$F = N E_o + m E_g + (1+m) E_{f,w} + W_e$$

onde $N$ (estoque original), $m$ (razão do gas‑cap/óleo) e outros parâmetros podem ser estimados por ajuste linear de $F$ versus os termos $E_o, E_g, E_{f,w}$.

\subsection{5.5 EBM para reservatórios de gás}
Para reservatórios dominados por gás a formulação é adaptada; utiliza‑se o fator de volume do gás $B_g$ e considera‑se o comportamento PVT do gás. Técnicas similares de linearização e plotagem (ex.: material balance gas plots) são usadas para estimar OGIP e responder por influxos.

\subsection{5.6 Técnicas de análise}
\begin{itemize}
\item Plotagens clássicas de material balance (ex.: $F$ vs $E_o$, $F$ vs combinação linear) para estimar $N$ e $m$.
\item Uso de regressão linear múltipla (Havlena & Odeh) para separar contribuições de compressibilidade e gas‑cap.
\item Análise de sensibilidade e propagação de incerteza (p.ex. Monte Carlo) para obter P10/P50/P90.
\end{itemize}

\subsection{5.7 Vantagens e limitações}
\begin{itemize}
\item Vantagens: não requer geometria detalhada; integra dados de produção; fornece estimativas robustas quando dados de produção e PVT são confiáveis.
\item Limitações: sensível a erros de PVT, incertezas em volumes injetados, e viés por heterogeneidade ou comunicação entre zonas; exige QA/QC rigoroso dos dados.
\end{itemize}

\subsection{5.8 Consolidação da aula — Exercícios}
\begin{enumerate}
\item Explique a lógica da linearização de Havlena & Odeh.
\item Dado um conjunto de dados de produção e PVT, descreva os passos para construir o termo $F$ e os termos $E_o,E_g,E_{f,w}$.
\item Discuta como incluir influxo de aquífero e injeção de água na formulação.
\end{enumerate}

\subsection{5.9 Tarefas}
\begin{enumerate}
\item Aplicar a linearização Havlena & Odeh a um conjunto de dados hipotético (fornecerei os dados se desejar).
\item Fazer uma análise de sensibilidade para $B_o$, $\phi$ e $S_{w,i}$ e discutir o impacto nas reservas estimadas.
\item Preparar um pequeno relatório comparando EBM e método volumétrico para um mesmo caso.
\end{enumerate}

\subsection{Definições importantes}
\begin{itemize}
\item **Termo $F$ (lado conhecido):** representa as quantidades medidas de produção/injeção devidamente convertidas para volumes de superfície usados na EBM.
\item **$N$ (estoque original):** volume originalmente em lugar de óleo (STB) ou gás (scf/Sm3).
\item **$m$ (razão gas‑cap/óleo):** relação entre volumes de gás livre no gas‑cap e volume de óleo no reservatório.
\item **$E_o, E_g, E_{f,w}$:** termos de variação volumétrica do óleo, do gás e do fluido de formação/água respectivamente (definidos na linearização Havlena & Odeh).
\end{itemize}

\subsection{Exemplo prático — aplicação simplificada da linearização Havlena & Odeh}
Considere um caso hipotético com os seguintes dados simplificados por período acumulado (valores ilustrativos):
\begin{itemize}
\item Produção acumulada de óleo $N_p = 100{,}000\,$STB;
\item Produção acumulada de gás $G_p = 50{,}000\,$scf (convertido para as mesmas unidades apropriadas);
\item $B_{o,i}=1.1$, $B_o$ atual = 1.15; $R_{s,i}=200\,$scf/STB, $R_s=180\,$scf/STB; $B_g=0.005$ respec.
\end{itemize}

Calcule um termo $E_o$ simplificado:
$$
E_o = B_o - B_{o,i} + (R_{s,i} - R_s)B_g
$$
Substituindo valores (exemplo):
$$
E_o = 1.15 - 1.10 + (200-180)\times0.005 = 0.05 + 20\times0.005 = 0.05 + 0.10 = 0.15
$$
Se repetirmos esse cálculo para várias datas e obtivermos $F$ correspondente, um ajuste linear de $F$ versus $E_o$ (junto com $E_g$ e $E_{f,w}$) permite estimar $N$ (coeficiente angular/intercepto dependendo da formulação).

\subsection{Observações práticas}
\begin{itemize}
\item Garantir que todos os volumes e fatores $B$ estejam convertidos para o mesmo sistema de unidades antes de montar $F$.
\item Corrigir por volumes injetados ($W_{inj},G_{inj}$) e por perdas/erros de medição.
\item Em presença de aquífero, incluir o termo $W_e$ de forma explícita e considerar modelos simplificados de influxo para estimativa.
\end{itemize}

\subsection{Glossário de símbolos (cap.5)}
\begin{itemize}
\item `F` — termo conhecido (volume produzido/injetado convertido para superfície).
\item `N` — estoque original de óleo (STB).
\item `m` — razão gas‑cap/óleo (adimensional).
\item `E_o, E_g, E_{f,w}` — termos de variação na linearização da EBM.
\item `B_o, B_g` — fatores de volume de formação do óleo e do gás.
\item `N_p, G_p` — produção acumulada de óleo e gás.
\item `W_e` — contribuição do influxo de água (aquífero).
\end{itemize}

\subsection{Dicas rápidas de estudo}
\begin{itemize}
\item Organize os dados de produção e PVT em uma planilha com colunas para $N_p,G_p,B_o,B_g,R_s$ e calcule $F,E_o,E_g,E_{f,w}$ periodicamente.
\item Plote pares $(F,E_o)$ e $(F,E_g)$ para observar linearidade e influências dominantes.
\item Teste sensibilidade a erros de PVT (p.ex. variação de $B_o$) para entender impacto nas estimativas de $N$.
\end{itemize}

---

Muito obrigado(a)!

% ---- file: exercícios_transcription.md



1



Engenharia de Reservatórios - I
Problemas e Soluções
Ano Lectivo 2023/2024
Docente
Geraldo Ramos, Bsc, Msc, PhD
Instituto Superior Politécnico de Tecnologias e Ciências - ISPTEC
18 de maio de 2024
COPYRIGHT: © 2024 by Geraldo Ramos. All rights reserved.
Nenhuma parte deste material pode ser produzido, armazenado, ou transmitido por qualquer meio electrô-
nico, mecânico, ou por fotocópia, gravação, digitalização, ou outros, sem prévia permissão por escrito do
autor.
Introdução
O objectivo deste documento é apresentar as soluções dos exercícios/problemas de cada aula da disciplina de
Engenharia de Reservatórios I, do 2◦ semestre, do 3◦ ano do curso de Engenharia de Petróleos do ISPTEC.
As respostas com a cor vermelha, ainda estão em desenvolvimento, excepto nas 
guras onde foi utilizada
para realçar informações ao passo que as de cor azul ou outra cor que não seja cor vermelha são consideradas
concluidas e veri
cadas.
Os estudantes são obrigados em resolver os exercícios antes de consultarem as respostas para o sucesso na
cadeira de Engenharia de Reservatórios I. Qualquer erro na solução agradecia noti
car o professor.
Os problemas deste documento estão relacionados com:
ˆ Capítulo 1  Introdução à Engenharia de Reservatórios;
ˆ Capítulo 2  Propriedades dos 
uidos;
ˆ Capítulo 3  Propriedade das rochas;
ˆ Capítulo 4  Integração de dados e Cálculo Volumétrico
ˆ Capítulo 5  Equação de balanço de materiais;
ˆ Capítulo 6 - Fluxo de Fluidos em Meios Porosos;
ˆ Capítulo 7 - Análise de Teste de Pressão Transiente em poços
ˆ Capítulo 8  In
uxo Natural de água;
ˆ Capítulo 9 - Análise de curvas de declínio e ajuste do histórico
O documento está dividido em três partes:
ˆ Parte - I: O questionário dos problemas
ˆ Parte - II: As soluções dos problemas
ˆ Bibliogra
a consultada
Espero que venha auxiliar para o seu sucesso nesta disciplina.
Geraldo Ramos, PhD
i

% ---- file: resumo_capitulos.md

\section{Resumo por Capítulo — Engenharia de Reservatórios I (Versão EXTREMAMENTE Detalhada)}

Este documento contém um resumo aprofundado dos capítulos 1–5, com fórmulas, definições, unidades, notas de interpretação e exemplos numéricos passo‑a‑passo. Use como referência de estudo e para gerar fichas de revisão.

Índice
\begin{itemize}
\item Capítulo 1 — Conceitos e sistema petrolífero
\item Capítulo 2 — Propriedades dos fluidos (PVT): teorias, fórmulas e correlações
\item Capítulo 3 — Propriedades das rochas: porosidade, permeabilidade, capilaridade e equações aplicadas
\item Capítulo 4 — Cálculo volumétrico (OOIP / OGIP): fórmulas de campo e SI, sensibilidade e Monte Carlo
\item Capítulo 5 — Equação de Balanço de Materiais (EBM): formulação, linearização e exemplos
\item Anexos: constantes, fatores de conversão, lista de fórmulas essenciais
\end{itemize}

---

\textbf{Nota de unidades:} Sempre indique o sistema: UNIDADES DE CAMPO (acres, ft, bbl, scf) ou SI (m, m³, Pa). Muitos fatores práticos (ex.: 7758) convertem acres·ft → bbl.

---

\subsection{Capítulo 1 — Conceitos e sistema petrolífero}

1.1 Principais definições
\begin{itemize}
\item Sistema petrolífero: conjunto de elementos necessários para geração, migração, armadilhamento e acumulação de hidrocarbonetos.
\item Sistema de produção: instalações e equipamentos para recuperação, elevação e tratamento dos fluidos do reservatório até a superfície.
\end{itemize}

1.2 Equações e relações fundamentais (úteis para este capítulo)
\begin{itemize}
\item Equilíbrio hidrostático (coluna incompressível ideal):
$$p(z)=p_{ref}+\rho g (Z_{ref}-Z)$$
onde $p$ em Pa (ou psi), $\rho$ densidade (kg/m³), $g$ = 9.80665 m/s², $Z$ profundidade.
\item Soma de saturações (misto trifásico):
$$S_w + S_o + S_g = 1$$
onde cada $S_\bullet$ é fração volumétrica adimensional.
\end{itemize}

1.3 Conceitos qualitativos importantes
\begin{itemize}
\item Armadilha (structural/estratigráfica/fraturada), rocha geradora/selante, migração; sincronismo e janela térmica (catagênese) — sem equações mas essenciais para interpretação geológica.
\end{itemize}

---

\subsection{Capítulo 2 — Propriedades dos fluidos (PVT)}

2.1 Variáveis termodinâmicas fundamentais
\begin{itemize}
\item Equação dos gases ideais (mol):
$$pV=nRT$$
onde $p$ (Pa), $V$ (m³), $n$ (mol), $R$ (8.314462618 J/(mol·K)), $T$ (K).
\item Fator de compressibilidade (Z):
$$Z=\dfrac{pV}{nRT}=\dfrac{p\,\overline{v}}{RT}$$
Z corrige o comportamento real do gás.
\end{itemize}

2.2 Definições PVT essenciais
\begin{itemize}
\item Fator de volume de formação do óleo ($B_o$):
$$B_o=\dfrac{V_{res\_oil}}{V_{surf\_oil}}\quad(\text{unidades: m}^3/\text{m}^3\;\text{ou bbl/STB})$$
\item Fator de volume do gás ($B_g$): volume de gás no reservatório por unidade de gás padrão (ft³/SCF ou m³/Sm³):
$$B_g=\dfrac{V_{res\_gas}}{V_{std\_gas}}$$
\item Razão gás/óleo dissolvido ($R_s$): normalmente em scf/STB (campo) ou m³/m³ (SI):
$$R_s=\dfrac{\text{volume de gás dissolvido no óleo (std)}}{\text{volume de óleo (superfície)}}$$
\end{itemize}

2.3 Compressibilidades e variações com pressão
\begin{itemize}
\item Compressibilidade do óleo (bo):
$$c_o = -\dfrac{1}{V_o}\dfrac{\mathrm{d}V_o}{\mathrm{d}p} = \dfrac{\mathrm{d}(\ln V_o)}{\mathrm{d}p}$$
\item Compressibilidade do gás (aproximação):
$$c_g \approx \dfrac{1}{p} \left(1 - \dfrac{\mathrm{d}\ln Z}{\mathrm{d}\ln p}\right)$$
Usar tabelas/curvas Z para cálculo prático.
\end{itemize}

2.4 Correlações e estimativas empíricas (nomes e uso)
\begin{itemize}
\item Correlações de $B_o$, $\mu_o$, $R_s$ (Standing, Vazquez‑Beggs, Beggs & Robinson para viscosidade). Estas fornecem estimativas quando não existem ensaios PVT.
\end{itemize}

2.5 Conversões úteis (campo ↔ SI)
\begin{itemize}
\item $1\,\text{scf}=0.0283168\,\text{m}^3$;
\item $1\,\text{STB}=0.1589873\,\text{m}^3$;
\item $1\,\text{bbl}=0.1589873\,\text{m}^3$.
\end{itemize}

2.6 Exemplo prático (passo a passo)
\begin{itemize}
\item Dados: $V_{res}=1.20\,$bbl, $V_{surf}=1.00\,$STB, $R_s=400\,$scf/STB.
\item $B_o = 1.20/1.00 = 1.20\,$bbl/STB.
\item $R_s$ em SI:
$$R_s(\text{m}^3/\text{m}^3)=\dfrac{400\times0.0283168}{0.1589873}\approx71.3\;\text{m}^3/\text{m}^3$$
\end{itemize}

---

\subsection{Capítulo 3 — Propriedades das rochas}

3.1 Porosidade e volumes
\begin{itemize}
\item Porosidade total:
$$\phi=\dfrac{V_p}{V_t}$$
onde $V_p$ é o volume de poros e $V_t$ o volume total.
\item Porosidade a partir de massa (método gravimétrico):
$$V_p=\dfrac{m_{sat}-m_{dry}}{\rho_f},\qquad \phi=\dfrac{V_p}{V_t}$$
\end{itemize}

3.2 Densidades e densidade aparente (logs)
\begin{itemize}
\item Massa específica/aparente: para rocha saturada
$$\rho_b=(1-\phi)\rho_{s}+\phi(S_w\rho_w+S_o\rho_o+S_g\rho_g)$$
onde $\rho_s$ é a densidade da matriz.
\end{itemize}

3.3 Permeabilidade (Lei de Darcy)
\begin{itemize}
\item Forma diferencial (unidimensional):
$$q=-\dfrac{kA}{\mu}\dfrac{\mathrm{d}p}{\mathrm{d}x}$$
onde $q$ é vazão volumétrica (m³/s), $k$ permeabilidade (m²), $A$ área (m²), $\mu$ viscosidade (Pa·s).
\item Fluxo radial estacionário (poço produtor, conduto linearizado):
$$q=\dfrac{2\pi k h (p_e-p_{wf})}{\mu\ln\left(\dfrac{r_e}{r_w}\right)}$$
onde $h$ é espessura produtiva, $r_e$ raio de drenagem, $r_w$ raio do poço.
\end{itemize}

3.4 Difusividade (equação do fluxo transiente)
\begin{itemize}
\item Para fluido ligeiramente compressível:
$$\dfrac{\partial p}{\partial t}=\dfrac{k}{\phi\mu c_t}\nabla^2 p$$
onde $c_t$ é compressibilidade total (soma das compressibilidades relevantes).
\end{itemize}

3.5 Capilaridade e função de Leverett
\begin{itemize}
\item Pressão capilar (tubo capilar simplificado):
$$P_c = \dfrac{2\sigma\cos\theta}{r}$$
\item Função de Leverett:
$$J(S_w)=\dfrac{P_c(S_w)\sqrt{k/\phi}}{\sigma\cos\theta}$$
\end{itemize}

3.6 Saturações e permeabilidades relativas
\begin{itemize}
\item $S_w+S_o+S_g=1$;
\item Permeabilidade relativa: $k_{r\alpha}=f(S_\alpha)$ (curvas experimentais).
\end{itemize}

3.7 Resistividade e saturação — Equação de Archie
\begin{itemize}
\item Forma usual (zona limpa, rocha não condutiva):
$$R_t = a\,R_w\,\phi^{-m}\,S_w^{-n}$$
ou invertendo para saturação:
$$S_w = \left(\dfrac{a\,R_w}{R_t\,\phi^{m}}\right)^{1/n}$$
com parâmetros empíricos $a,m,n$.
\end{itemize}

3.8 Exemplos resolvidos (porosidade por massa; Arquímedes)
\begin{itemize}
\item Ver exemplos do capítulo 3 (por ex., $\phi\approx16{,}5\%$ e $22{,}2\%$ em exercícios práticos — seguir passos mostrados).
\end{itemize}

---

\subsection{Capítulo 4 — Integração de Dados e Cálculo Volumétrico}

4.1 Forma geral (reservoir and surface)
\begin{itemize}
\item Volume de fluido em reservatório (m³):
$$V_{fluid,res}=V_r\,\phi\,(1-S_w)$$
\item Conversão para superfície usando fator de volume $B$ (m³/m³):
$$N_{surface}=\dfrac{V_{fluid,res}}{B}$$
\end{itemize}

4.2 Fórmulas práticas (unidades de campo)
\begin{itemize}
\item OOIP (STB):
$$OOIP=\dfrac{7758\,A\,h\,\phi\,(1-S_w)}{B_o}$$
onde $A$ em acres, $h$ em ft, $\phi$ fração, $B_o$ em bbl/STB.\
Derivação rápida do fator 7758:
1 acre = 43560 ft^2; 1 acre·ft = 43560 ft^3; 1 bbl = 5.614583 ft^3 → 43560/5.614583 \approx 7758 bbl/acre·ft.
\end{itemize}

\begin{itemize}
\item OGIP (SCF):
$$OGIP=\dfrac{43560\,A\,h\,\phi\,(1-S_w)}{B_g}$$
onde $B_g$ é o fator volume do gás (ft^3/SCF).
\end{itemize}

4.3 Estimativa de reservas recuperáveis
\begin{itemize}
\item Reservas volumétricas (por fator de recuperação $F_R$):
$$N_R = OOIP\times F_R$$
ou pontualmente (contando saturações residuais):
$$N_R = V_r\,\phi\,(1-S_{w,i})\left(\dfrac{S_{o,i}}{B_{o,i}} - \dfrac{S_{o,r}}{B_{o,r}}\right)$$
\end{itemize}

4.4 Análise de sensibilidade e incerteza
\begin{itemize}
\item Derivada parcial (sensibilidade de OOIP a $\phi$):
$$\dfrac{\partial OOIP}{\partial\phi}=\dfrac{7758\,A\,h\,(1-S_w)}{B_o}$$
\item Monte Carlo: descreva distribuições para $A,h,\phi,S_w,B_o$ (p.ex. triangular, normal truncada, lognormal), gere N amostras e calcule percentis P10/P50/P90 da distribuição resultante de OOIP.
\end{itemize}

4.5 Boas práticas de integração
\begin{itemize}
\item QA/QC em dados de core, logs e PVT; uso de mapas de net pay e cut‑offs; segmentação setorial (zonas) — calcular OOIP por setor e somar.
\end{itemize}

---

\subsection{Capítulo 5 — Equação de Balanço de Materiais (EBM)}

5.1 Conceito geral
\begin{itemize}
\item A EBM relaciona as quantidades produzidas/injetadas com a variação de estoque no reservatório; é uma equação de conservação de massa para os poros.
\end{itemize}

5.2 Forma linearizada (Havlena \& Odeh)
\begin{itemize}
\item Variável conhecida (lado esquerdo), $F$, é construída a partir de dados de produção e injeção convertidos para volumes de superfície. A forma linearizada usada comumente é:
$$F = N\,E_o + m\,E_g + (1+m)\,E_{f,w} + W_e$$
onde:
\item $E_o = B_o - B_{o,i} + (R_{s,i}-R_s)\,B_g$;
\item $E_g = B_{o,i}\left(\dfrac{B_g}{B_{g,i}} - 1\right)$;
\item $E_{f,w} = B_{o,i}\left(S_{w,i}c_w + \dfrac{c_f}{1-S_{w,i}}\Delta p\right)$;
\item $N$ é OOIP (a estimar), $m$ razão gas‑cap/óleo (a estimar), $W_e$ contribuição líquida do aquífero.
\end{itemize}

5.3 Construção prática de $F$ (exemplo simplificado)
\begin{itemize}
\item Exemplo (formas ilustrativas): calcular $F$ acumulado por intervalo como combinação de termos de produção $N_p B_o$, $G_p$ e correções por $R_s$, volumes injetados, etc. (ver exercícios e ficheiro cap5 para forma usada no curso).
\end{itemize}

5.4 Estimação por regressão
\begin{itemize}
\item Monte as colunas $F, E_o, E_g, E_{f,w}$ para vários instantes. Execute regressão linear múltipla:
$$F = \beta_1 E_o + \beta_2 E_g + \beta_3 E_{f,w} + \varepsilon$$
onde idealmente $\beta_1=N$, $\beta_2=m$, $\beta_3=(1+m)$ e $\varepsilon$ residuais (ou ajustar conforme convenção usada).
\end{itemize}

5.5 Notas sobre gás e p/z
\begin{itemize}
\item Para reservatórios gas‑dominados, a análise p/z (plot de $p/z$ vs $G_p$) permite estimar OGIP por extrapolação; p/z é inversamente proporcional ao volume remanescente quando as condições são adequadas. (Implementação numérica e tratamentos de retroinjeção exigem cuidados de unidades e correções).
\end{itemize}

---

\subsection{Anexos — Fórmulas essenciais e constantes}

Constantes e conversões rápidas
\begin{itemize}
\item $1\,\text{acre}=43560\,\text{ft}^2$.
\item $1\,\text{acre·ft}=43560\,\text{ft}^3\approx7758\,\text{bbl}$.
\item $1\,\text{bbl}=0.1589873\,\text{m}^3$.
\item $1\,\text{scf}=0.0283168\,\text{m}^3$.
\item $1\,\text{D}=9.869233\times10^{-13}\,\text{m}^2$.
\item $1\,\text{psi}=6894.757\,\text{Pa}$.
\end{itemize}

Lista condensada de fórmulas (referência rápida)
\begin{itemize}
\item Somatório das saturações: $S_w+S_o+S_g=1$.
\item API gravity: $API=\dfrac{141.5}{SG}-131.5$.
\item Darcy (unidimensional): $q=-\dfrac{kA}{\mu}\dfrac{\mathrm{d}p}{\mathrm{d}x}$.
\item Radial steady flow to well: $q=\dfrac{2\pi k h (p_e-p_{wf})}{\mu\ln(r_e/r_w)}$.
\item OOIP (campo): $OOIP=\dfrac{7758Ah\phi(1-S_w)}{B_o}$.
\item OGIP (campo): $OGIP=\dfrac{43560Ah\phi(1-S_w)}{B_g}$.
\item Porosidade: $\phi=V_p/V_t$.
\item Compressibilidade: $c= -\dfrac{1}{V}\dfrac{\mathrm{d}V}{\mathrm{d}p}$.
\item Archie: $R_t=aR_w\phi^{-m}S_w^{-n}$.
\item Capilaridade (tubo): $P_c=\dfrac{2\sigma\cos\theta}{r}$.
\item Leverett: $J(S_w)=\dfrac{P_c(S_w)\sqrt{k/\phi}}{\sigma\cos\theta}$.
\end{itemize}

---

\subsection{Sugestões de estudo e uso deste ficheiro}
\begin{itemize}
\item Transforme cada subseção num cartão de revisão (Anki) — fórmulas, definições e exemplos numéricos.
\item Implemente as fórmulas essenciais numa planilha (OOIP/OGIP/Monte Carlo) para treino prático.
\item Para a prova: foque em interpretar o significado físico das fórmulas (o que altera OOIP, por que $B_o$ varia com pressão, como $R_s$ altera mobilidade, etc.).
\end{itemize}

---

Arquivo gerado automaticamente pelo assistente — versão detalhada criada a pedido do utilizador.

% ---- file: resumo_capitulos_ultra_detalhado.md
\section{Resumo por Capítulo — Engenharia de Reservatórios I (VERSÃO ULTRA‑DETALHADA)}

Este ficheiro contém uma compilação exaustiva de fórmulas, definições, derivações e exemplos numéricos para os capítulos 1–5. Use como referência técnica aprofundada durante a preparação.

Índice
\begin{itemize}
\item Capítulo 1 — Sistema petrolífero e conceitos fundamentais
\item Capítulo 2 — Propriedades dos fluidos (PVT): equações, correlações e cálculos
\item Capítulo 3 — Propriedades de rochas: porosidade, permeabilidade, capilaridade, testes
\item Capítulo 4 — Cálculo volumétrico (OOIP / OGIP): fórmulas de campo e SI, sensibilidade, Monte Carlo
\item Capítulo 5 — Equação de Balanço de Materiais (EBM): formulação, linearização, p/z e exemplos
\item Anexos: constantes, fatores de conversão, valores típicos, checklist de passos práticos
\end{itemize}

---

\textbf{Observação sobre unidades}
\begin{itemize}
\item Indique sempre o sistema: CAMPO (acres, ft, bbl, scf, psi) ou SI (m, m³, Pa). Mantenha consistência.
\item Fatores de conversão essenciais estão no Anexo.
\end{itemize}

---

\subsection{Capítulo 1 — Sistema petrolífero e conceitos fundamentais}

\begin{enumerate}
\item Definições essenciais
\begin{itemize}
\item Sistema petrolífero: rocha geradora + rocha selante + armadilha + migração + sincronismo. 
\item Sistema de produção: conjunto de poços, tubulações, elevação artificial, separadores, tratamento e escoamento.
\end{itemize}
\end{enumerate}

\begin{enumerate}
\item Relações e equações úteis
\begin{itemize}
\item Equilíbrio hidrostático (coluna vertical):
$$p(Z)=p_{ref}+\int_{Z_{ref}}^{Z}\rho(z') g\,dz'\approx p_{ref}+\rho g (Z_{ref}-Z)$$
(usar densidade local em cada camada quando heterogênea)
\end{itemize}
\end{enumerate}

\begin{itemize}
\item Soma das saturações (para mistura trifásica):
$$S_w+S_o+S_g=1$$
\end{itemize}

\begin{enumerate}
\item Notas interpretativas
\begin{itemize}
\item Identificação do tipo de reservatório (oil‑drive, gas‑cap, water‑drive, solution gas drive) depende de presença de gas cap, aquífero e relação GOR/Rs.
\end{itemize}
\end{enumerate}

---

\subsection{Capítulo 2 — Propriedades dos fluidos (PVT)}

2.1. Grandezas fundamentais
\begin{itemize}
\item Equação dos gases (mol):
$$pV=nRT$$
\item Fator de compressibilidade (Z):
$$Z=\dfrac{pV}{nRT}$$
\end{itemize}

2.2. Fatores de volume e razões
\begin{itemize}
\item Fator de volume de formação do óleo (óleo: reservatório → superfície):
$$B_o=\dfrac{V_{res\_oil}}{V_{surf\_oil}}\quad(\text{m}^3/\text{m}^3\;\text{ou bbl/STB})$$
\item Relação densidade ↔ B_o: massa invariável entre estados
$$\rho_{res} = \dfrac{m}{V_{res}},\qquad \rho_{surf}=\dfrac{m}{V_{surf}}\Rightarrow B_o=\dfrac{V_{res}}{V_{surf}}=\dfrac{\rho_{surf}}{\rho_{res}}$$
\end{itemize}

\begin{itemize}
\item Razão gás/óleo dissolvido:
$$R_s=\dfrac{\text{vol. gás liberado (condições padrão)}}{\text{vol. óleo (superfície)}}\quad(\text{scf/STB ou m}^3/\text{m}^3)$$
\end{itemize}

2.3. Compressibilidade de fluidos
\begin{itemize}
\item Compressibilidade do óleo (adimensional, psi⁻1 ou Pa⁻1):
$$c_o = -\dfrac{1}{V_o}\dfrac{\mathrm{d}V_o}{\mathrm{d}p}=\dfrac{\mathrm{d}(\ln V_o)}{\mathrm{d}p}$$
\item Compressibilidade do gás (aprox.):
$$c_g \approx \dfrac{1}{p}\left(1-\dfrac{\partial\ln Z}{\partial\ln p}\right)$$
(obter Z(p,T) de charts ou EoS)
\end{itemize}

2.4. Equações de estado cúbicas (uso em PVT)
\begin{itemize}
\item Definições gerais (PR e SRK usadas com frequência):
\item Definir constantes críticas de cada componente: $T_c, p_c, \omega$ (fator acêntrico).
\end{itemize}

\begin{itemize}
\item Peng‑Robinson (PR):
\item Parâmetros puros:
  $$a = 0.45724\dfrac{R^2 T_c^2}{p_c},\qquad b = 0.07780\dfrac{R T_c}{p_c}$$
\item Fator de temperatura:
  $$\alpha(T)=\left[1+\kappa (1-\sqrt{T_r})\right]^2,\quad T_r=\dfrac{T}{T_c}$$
  com
  $$\kappa = 0.37464 + 1.54226\omega - 0.26992\omega^2.$$
\item EoS:
  $$p=\dfrac{RT}{V-b}-\dfrac{a\alpha(T)}{V(V+b)+b(V-b)}$$
\end{itemize}

\begin{itemize}
\item Soave‑Redlich‑Kwong (SRK):
\item Parâmetros:
  $$a = 0.42748\dfrac{R^2T_c^2}{p_c},\qquad b=0.08664\dfrac{RT_c}{p_c}$$
\item EoS (forma):
  $$p=\dfrac{RT}{V-b}-\dfrac{a\alpha(T)}{V(V+b)}$$
\item Kappa (SRK): parâmetro dependente de $\omega$ (ver formulação SRK).
\end{itemize}

\begin{itemize}
\item Mistura (regra de mistura van der Waals tipo):
  $$a_{mix}=\sum_i\sum_j x_i x_j \sqrt{a_i a_j}(1-k_{ij}),\qquad b_{mix}=\sum_i x_i b_i$$
  onde $k_{ij}$ são parâmetros de interação binária.
\end{itemize}

\begin{itemize}
\item Redução a polinomial em Z (cúbica): em geral obtem‑se um polinómio cúbico em $Z$ (derivado do EoS multiplicado por $V$ e normalizado) do tipo:
  $$Z^3 + c_2 Z^2 + c_1 Z + c_0 = 0$$
  com coeficientes que dependem de $A$ e $B$:
  $$A=\dfrac{a p}{(R T)^2},\qquad B=\dfrac{b p}{R T}$$
  e um polinómio padrão (usado tanto para PR como SRK após definição apropriada de $A,B$):
  $$Z^3-(1-B)Z^2+(A-3B^2-2B)Z-(AB-B^2-B^3)=0$$
  (resolver numericamente para raízes reais; raízes correspondem a fases gasosa/líquida onde aplicável).
\end{itemize}

2.5. Cálculo de fatores de volume com Z
\begin{itemize}
\item Formação volume factor do gás (por molar/molar base):
  $$B_g = \dfrac{Z R T}{p}\cdot\dfrac{V_{ref\_units}}{n_{ref}}$$
  (usar forma prática/constantes de conversão ao trabalhar em scf/ft³ ou SI).
\end{itemize}

2.6. Correlações empíricas e conversões úteis
\begin{itemize}
\item Conversões:
  $$1\,\text{scf}=0.0283168\,\text{m}^3,\quad1\,\text{STB}=0.1589873\,\text{m}^3$$
\item API gravity (relativa à água a 60°F):
  $$API=\dfrac{141.5}{SG_{60°F}} -131.5,\quad SG=\dfrac{\rho_{oil}}{\rho_{water}}$$
\end{itemize}

2.7. Exemplo prático PVT (completo)
\begin{itemize}
\item Dados: $V_{res}=1.20\,$bbl; $V_{surf}=1.00\,$STB; $R_s=400\,$scf/STB.
\begin{enumerate}
\item $B_o=1.20/1.00=1.20\,$bbl/STB.
\item $R_s(\text{SI})=400\times0.0283168/0.1589873\approx71.3\,\text{m}^3/\text{m}^3$.
\item Densidade de reservatório: $\rho_{res}=\rho_{surf}/B_o$.
\end{itemize}
\end{enumerate}

---

\subsection{Capítulo 3 — Propriedades das rochas (detalhado)}

3.1. Porosidade (definições e métodos)
\begin{itemize}
\item Porosidade total: $\phi=V_p/V_t$.
\item Métodos medição: gravimétricos (pesagem seca/saturada), porosimetria de mercúrio, NMR, micro‑CT.
\item Porosidade efetiva (conectada) e irreducível (connate water) — distinguir para flow modelling.
\end{itemize}

3.2. Densidade de rocha saturada (mixing law)
$$\rho_b = (1-\phi)\rho_s + \phi (S_w \rho_w + S_o \rho_o + S_g \rho_g)$$
onde $\rho_s$ densidade da matriz.

3.3. Permeabilidade e Lei de Darcy
\begin{itemize}
\item Forma integral (unidimensional):
$$q = -\dfrac{k A}{\mu}\dfrac{\Delta p}{L}$$
\item Para fluxo radial permanente para poço produtor (campo/ft):
$$q=\dfrac{2\pi k h (p_e-p_{wf})}{\mu \ln(r_e/r_w)}$$
(usar conversões quando q em STB/d, µ em cP, k em mD — ver fórmulas de campo no anexo)
\end{itemize}

3.4. Equação de difusividade (transiente, ligeiramente compressível)
\begin{itemize}
\item Forma geral:
$$\dfrac{\partial p}{\partial t}=\dfrac{k}{\mu S} \nabla^2 p$$
onde $S$ (coeficiente de armazenamento) geralmente
$$S=\phi c_f + (1-\phi)c_s$$
com $c_f$ compressibilidade do fluido e $c_s$ compressibilidade da matriz sólida.
\end{itemize}

3.5. Solução transiente radial (analítica aproximada — semilog)
\begin{itemize}
\item Para regime semilog (pseudosteady), pressão medida no poço varia com tempo t:
$$p(r_w,t)=p_i - \dfrac{q B \mu}{4\pi k h}\left[\ln\left(\dfrac{4 k t}{\phi \mu c_t r_w^2}\right)-0.80907\right]$$
\item Em base‑10 logs (análise semilog): a declividade $m$ (psi por ciclo log10) é
$$m=\dfrac{2.303\,q B \mu}{4\pi k h} = \dfrac{0.183 q B \mu}{k h}$$
  e, isolando $k$ (unidades de campo com $q$ em STB/d, $\mu$ cP, $h$ ft, $m$ psi/log10):
$$k(\text{mD}) = \dfrac{162.6\,q\,B\,\mu}{h\,m}$$
(162.6 é constante empírica para conversões entre unidades; verificar convenções de unidades)
\end{itemize}

3.6. Função de Leverett e pressão capilar
\begin{itemize}
\item Tensão interfacial e rádio capilar simplificado (tubo):
$$P_c=\dfrac{2\sigma \cos\theta}{r}$$
\item Função J de Leverett:
$$J(S_w)=\dfrac{P_c(S_w)\sqrt{k/\phi}}{\sigma\cos\theta}$$
(usar J(S) para escalonar curvas Pc entre rochas com diferentes k e φ)
\end{itemize}

3.7. Permeabilidade relativa \& modelos
\begin{itemize}
\item Modelo de Corey (exemplo):
$$S_{we}=\dfrac{S_w-S_{wr}}{1-S_{or}-S_{wr}}$$
$$k_{rw}=k_{rw0} S_{we}^{n_w},\quad k_{ro}=k_{ro0}(1-S_{we})^{n_o}$$
\item Parâmetros típicos: $n_w, n_o$ entre 2–4 dependendo de rocha/molhabilidade.
\end{itemize}

3.8. Resistividade e Archie
\begin{itemize}
\item Equação de Archie:
$$R_t = a R_w \phi^{-m} S_w^{-n}$$
\item Parâmetros empíricos: $a\approx1$, $m\approx1.8-2.2$, $n\approx2$ (variam com litologia)
\end{itemize}

---

\subsection{Capítulo 4 — Cálculo volumétrico (OOIP / OGIP) — detalhe completo}

4.1. Fórmula fundamental (reservoir → superfície)
\begin{itemize}
\item Volume de fluido em reservatório (m³):
$$V_{fluid,res}=V_r\,\phi\,(1-S_w)$$
\item Volume de superfície:
$$N_{surface}=\dfrac{V_{fluid,res}}{B}\quad(\text{B = fator de volume de formação})$$
\end{itemize}

4.2. Fórmulas práticas de campo
\begin{itemize}
\item OOIP (barris stock‑tank):
$$OOIP=\dfrac{7758\,A\,h\,\phi\,(1-S_w)}{B_o}$$
\item Derivação do factor 7758:
\item 1 acre = 43,560 ft²; 1 ft³ = 0.178107... bbl? (usar conversões exatas). A relação completa é: 43560 ft² × 1 ft = 43560 ft³ por acre·ft; 1 bbl = 5.614583 ft³; 43560/5.614583 ≈ 7758 bbl/acre·ft.
\end{itemize}

\begin{itemize}
\item OGIP (scf):
$$OGIP=\dfrac{43560\,A\,h\,\phi\,(1-S_w)}{B_g}$$
(43560 = ft³/acre coeficient)
\end{itemize}

4.3. Sensibilidade analítica (derivadas parciais)
\begin{itemize}
\item Sensibilidade de OOIP a φ:
$$\dfrac{\partial OOIP}{\partial\phi}=\dfrac{7758 A h (1-S_w)}{B_o}$$
\item Sensibilidade fraccional relativa: usar variáveis aleatórias e Monte Carlo para propagar incertezas.
\end{itemize}

4.4. Monte Carlo (passos práticos)
\begin{enumerate}
\item Definir distribuições para cada variável incerta (A,h,φ,S_w,B_o): triangular, normal truncada ou lognormal conforme justificativa.
\item Amostrar N vezes (p.ex. 10k‑100k iterações para robustez) e calcular OOIP para cada amostra.
\item Ordenar resultados e extrair percentis P10/P50/P90.
\item Reportar média, mediana (P50), intervalo de confiança e curvas CDF/PDF.
\end{enumerate}

4.5. Estimativa setorial
\begin{itemize}
\item Dividir reservatório em setores i com A_i,h_i,φ_i,S_{w,i} e somar:
$$OOIP_{total}=\sum_i \dfrac{7758 A_i h_i \phi_i (1-S_{w,i})}{B_{o,i}}$$
\end{itemize}

---

\subsection{Capítulo 5 — Equação de Balanço de Materiais (EBM) — completo}

5.1. Conceito (conservação de massa)
\begin{itemize}
\item A EBM parte do princípio: estoque inicial = estoque atual + produção acumulada - injeções + influxo. Convertendo volumes e considerando compressibilidade resulta em equações relacionando variação de pressão com quantidades produzidas/injetadas.
\end{itemize}

5.2. Forma linearizada (Havlena \& Odeh, 1963)
\begin{itemize}
\item Expressão usada em análise prática:
$$F = N E_o + m E_g + (1+m) E_{f,w} + W_e$$
onde:
\item $F$ termo conhecido (construído a partir de produção/injeção e PVT);
\item $N$ estoque original de óleo (OOIP);
\item $m$ razão gas‑cap/óleo (adimensional);
\item $E_o,E_g,E_{f,w}$ termos calculáveis a partir de PVT/pressão;
\item $W_e$ termo de influxo do aquífero.
\end{itemize}

5.3. Definições práticas (forma operacional)
\begin{itemize}
\item Termos (formas comumente usadas em regressão):
$$E_o = B_o - B_{o,i} + (R_{s,i}-R_s)B_g$$
$$E_g = B_{o,i}\left(\dfrac{B_g}{B_{g,i}} - 1\right)$$
$$E_{f,w}=B_{o,i}\left(S_{w,i}c_w + \dfrac{c_f}{1-S_{w,i}}\Delta p\right)$$
\item Lado conhecido $F$ (exemplo):
$$F = N_p B_o + G_p - R_s B_g + W_p B_w - W_{inj} B_w - G_{inj} B_{g,inj}$$
(ajustar sinais e convenções conforme o conjunto de dados — ver notas do curso)
\end{itemize}

5.4. Procedimento prático
\begin{enumerate}
\item Organizar série temporal: p(t), N_p(t), G_p(t), volumes injetados W_{inj},G_{inj}.
\item Calcular $B_o(p),B_g(p),R_s(p)$ a partir de PVT ou correlações.
\item Calcular colunas $E_o,E_g,E_{f,w}$ para cada instante.
\item Calcular $F$ para cada instante.
\item Executar regressão linear múltipla: ajustar $F$ por combinação linear dos termos para obter $N$ e $m$ (coeficientes de regressão) e estimar $W_e$.
\end{enumerate}

5.5. Método p/z (para gás)
\begin{itemize}
\item Plot de $p/z$ vs $G_p$: em reservatórios gas‑dominados volumétricos sem influxo, a extrapolação linear pode fornecer OGIP. Interpretação: declive e intercepto do plot relacionam OGIP e condições iniciais; trate com cuidado correções de temperatura e compressibilidade.
\end{itemize}

5.6. Exemplo simplificado (numérico)
\begin{itemize}
\item Construir tabela com p, N_p, G_p, B_o(p),B_g(p),R_s(p), calcular $E$'s e ajustar. (ver tarefas do capítulo 5 para exemplo completo do curso).
\end{itemize}

---

\subsection{Anexos — constantes, fatores de conversão e valores típicos}

Constantes importantes:
\begin{itemize}
\item $R_{universal}=8.314462618\;\text{J/(mol·K)}$ (usar unidades coerentes)
\item $g=9.80665\;\text{m/s}^2$
\end{itemize}

Conversões úteis:
\begin{itemize}
\item $1\,\text{acre}=43560\,\text{ft}^2$
\item $1\,\text{acre·ft}=43560\,\text{ft}^3$
\item $1\,\text{bbl}=0.1589873\,\text{m}^3$
\item $1\,\text{scf}=0.0283168\,\text{m}^3$
\item $1\,\text{D}=9.869233\times10^{-13}\,\text{m}^2$
\item $1\,\text{psi}=6894.757\,\text{Pa}$
\end{itemize}

Valores típicos (ordem de grandeza)
\begin{itemize}
\item Porosidade: 5%–30% para rochas reservatório; médias úteis 10%–25%.
\item Permeabilidade: mD–D (argiloso <1 mD, arenito bom 100–1000 mD, carbonatos variam muito).
\item Exponentes de Archie: $m\approx1.8-2.2$, $n\approx2$.
\end{itemize}

Checklist prático para cada capítulo (sintético)
\begin{itemize}
\item Cap.1: listar elementos do sistema petrolífero; identificar tipo de armadilha e presença de gas‑cap/aqüífero.
\item Cap.2: ter curvas B_o(p), R_s(p), µ_o(p) e tabela PVT; conhecer EoS e quando aplicá‑las.
\item Cap.3: consolidar porosidade/permeabilidade por core e logs; obter curvas Pc(S) e kr(S).
\item Cap.4: montar planilha volumétrica setorial; rodar sensibilidade e Monte Carlo.
\item Cap.5: organizar séries de produção; calcular colunas E; aplicar regressão e interpretar resultados.
\end{itemize}

---

FIM — ficheiro gerado como versão ULTRA‑DETALHADA. Se desejar, posso:
\begin{itemize}
\item (A) sobrepor `resumo_capitulos.md` com esta versão;
\item (B) gerar PDF/LaTeX desta versão; ou
\item (C) extrair cartões Anki (Q/A) automaticamente a partir das definições e fórmulas.
\end{itemize}

\subsection{Exercícios resolvidos (seleção representativa)}

Esta secção apresenta soluções passo‑a‑passo para exercícios-chave dos capítulos 1–5. Use estes exemplos como modelo para resolver problemas semelhantes.

\subsubsection{Capítulo 1 — Verdadeiro/Falso (respostas e justificativas breves)}
\begin{enumerate}
\item V — Armadilha geológica (structural/estratigráfica) impede migração e acumula hidrocarbonetos.
\item F — Rocha geradora é rica em matéria orgânica (não baixa).
\item F — Rocha selante tem **baixa** permeabilidade e impede fluxo.
\item V — Migração ocorre por meios porosos e fraturados, favorável a caminhos permeáveis.
\item F — Sincronismo é relevante: geração, migração e armadilhamento devem coincidir temporalmente.
\item V — Catagênese = craqueamento térmico do querogénio em hidrocarbonetos.
\item V — Sistema de produção inclui coleta, elevação e separação até a superfície.
\item V — Em gas‑cap drive o gás livre expande e ajuda a manter pressão.
\item F — Porosidade e permeabilidade são propriedades distintas (volume de poros vs facilidade de fluxo).
\item V — $B_o$ relaciona volumes em reservatório e superfície (reservoir → surface).
\end{enumerate}

\subsubsection{Capítulo 2 — PVT: exemplo resolvido (cálculo de $B_o$ e conversão de $R_s$)}
Enunciado: Num ensaio PVT por 1 STB de óleo obteve‑se $V_{res}=1.20\,$bbl e $V_{surf}=1.00\,$STB; gás libertado $R_s=400\,$scf/STB.

\begin{enumerate}[label=\Alph*)]
\item Cálculo de $B_o$:\n$$B_o=\dfrac{V_{res}}{V_{surf}}=\dfrac{1.20}{1.00}=1.20\;\text{bbl/STB}.$$ 
\end{enumerate}

\begin{enumerate}[label=\Alph*)]
\item Conversão de $R_s$ para m³/m³:\n1 scf = 0.0283168 m³; 1 STB = 0.1589873 m³.
$$R_s(\text{m}^3/\text{m}^3)=\dfrac{400\times0.0283168}{0.1589873}\approx\dfrac{11.32672}{0.1589873}\approx71.3\;\text{m}^3/\text{m}^3.$$ 
\end{enumerate}

\begin{enumerate}[label=\Alph*)]
\item Interpretação breve: $B_o>1$ indica que o volume no reservatório é maior que o volume final de superfície por unidade (efeitos de compressibilidade e gás dissolvido). $R_s$ alto implica presença significativa de gás dissolvido — ao atingir o ponto de bolha o gás liberta‑se, alterando mobilidade.\n\end{enumerate}

\subsubsection{Capítulo 3 — Rochas: exercícios resolvidos (porosidade por massa e método de Arquímedes)}

**Exemplo (Exercício 1)** — Dados: $m_{sat}=130\,$g; $m_{dry}=105\,$g; $\rho_o=0.84\,$g/cm^3; $V_t=180\,$cm^3.

\begin{enumerate}
\item Volume de fluido nos poros:
$$V_f=\dfrac{m_{sat}-m_{dry}}{\rho_o}=\dfrac{130-105}{0.84}=\dfrac{25}{0.84}\approx29.7619\;\text{cm}^3.$$ 
\end{enumerate}

\begin{enumerate}
\item Porosidade:
$$\phi=\dfrac{V_p}{V_t}=\dfrac{29.7619}{180}\approx0.16534\approx16.53\%.$$ 
\end{enumerate}

**Exemplo (Exercício 5 — método de Arquímedes)** — Dados: $m_{dry}=330\,$g; $m_{sat}=360\,$g; $m_{ap\_agua}=225\,$g; $\rho_{agua}=1\,$g/cm^3.

\begin{enumerate}
\item Volume total da amostra via empuxo (diferença entre peso em ar e peso aparente em água):
$$V_t=\dfrac{m_{sat}-m_{ap\_agua}}{\rho_{agua}}=\dfrac{360-225}{1}=135\;\text{cm}^3.$$ 
\end{enumerate}

\begin{enumerate}
\item Volume poroso (volume de fluido nos poros):
$$V_p=m_{sat}-m_{dry}=360-330=30\;\text{cm}^3.$$ 
\end{enumerate}

\begin{enumerate}
\item Porosidade:
$$\phi=\dfrac{V_p}{V_t}=\dfrac{30}{135}\approx0.22222\approx22.22\%. $$
\end{enumerate}

Observação: nos relatórios, apresente as unidades, arredondamento e possíveis fontes de erro experimental.

\subsubsection{Capítulo 4 — Cálculo volumétrico: exemplo resolvido (OOIP)}

Dados (exercício 4.8): $A=200\,$acres; $h_{net}=30\,$ft; $\phi=0.18$; $S_{wi}=0.25$; $B_o=1.2\,$bbl/STB.

Fórmula prática:
$$OOIP=\dfrac{7758\,A\,h\,\phi\,(1-S_w)}{B_o}.$$ 

Substituindo os valores e calculando passo a passo:
\begin{align*}
7758\times200&=1\,551\,600\\
1\,551\,600\times30&=46\,548\,000\\
46\,548\,000\times0.18&=8\,378\,640\\
8\,378\,640\times0.75&=6\,283\,980\\
OOIP&=\dfrac{6\,283\,980}{1.2}\approx5\,236\,650\;\text{STB}.
\end{align*}

Interpretação: aproximadamente 5.24×10^6 STB originalmente em lugar.

\subsubsection{Capítulo 5 — EBM: exemplo prático (cálculo de $E_o$ e montagem de colunas)}

Dados ilustrativos (simplificados): $B_{o,i}=1.10$, $B_o=1.15$, $R_{s,i}=200\,$scf/STB, $R_s=180\,$scf/STB, $B_g=0.005$ (unidades consistentes com a formulação do curso).

Cálculo do termo $E_o$ (Havlena \& Odeh):
$$E_o = B_o - B_{o,i} + (R_{s,i}-R_s)B_g$$
Substituindo:
$$E_o = 1.15 - 1.10 + (200-180)\times0.005 = 0.05 + 20\times0.005 = 0.05 + 0.10 = 0.15.$$ 

Notas práticas para montar a tabela de análise EBM:
\begin{itemize}
\item Para cada instante (data) calcule: $p$, $N_p$, $G_p$, $B_o(p)$, $B_g(p)$, $R_s(p)$.
\item Calcule colunas $E_o,E_g,E_{f,w}$ por fórmulas definidas; calcule $F$ (lado conhecido) usando volumes produzidos/injetados convertidos para unidades compatíveis;
\item Execute regressão linear múltipla de $F$ versus colunas $E_o,E_g,E_{f,w}$ para obter estimativas de $N$ (coeficiente associado a $E_o$) e $m$ (coeficiente associado a $E_g$).
\end{itemize}

Exemplo de output parcial (apenas ilustrativo):
| t | p (psi) | N_p (STB) | G_p (scf) | B_o | B_g | R_s | E_o | E_g | E_{f,w} | F |
|---:|---:|---:|---:|---:|---:|---:|---:|---:|---:|---:|
| t_1 | 3000 | 100000 | 50000 | 1.10 | 0.0050 | 200 | 0.00 | 0.00 | 0.00 | F_1 |
| t_2 | 2900 | 120000 | 60000 | 1.12 | 0.0051 | 195 | 0.02 | 0.01 | 0.002 | F_2 |

Onde $F_i$ é montado conforme convensão adotada no curso; preste atenção às unidades (STB vs scf) e às conversões necessárias.

---

Se desejar, posso expandir esta secção com soluções de todos os exercícios presentes em \texttt{exercícios_transcription.md} (isto exigirá mais tempo e poderei criar um ficheiro separado \texttt{exercicios_resolvidos.md}).

Arquivo criado automaticamente pelo assistente.

% ---- file: compressibilidade_Z_cheatsheet.md
\section{Compressibilidade (fator) Z — Cheatsheet extremo}

Objetivo: referência completa e passo‑a‑passo para entender, calcular e aplicar o fator de compressibilidade $Z$ (gás real), com métodos práticos (Standing–Katz, EoS), derivadas úteis e um exemplo numérico com código Python.

---

\textbf{Resumo rápido}
\begin{itemize}
\item **Definição:** $Z = \dfrac{p V_m}{R T}$, mede o desvio do gás em relação ao gás ideal ($Z=1$ ideal).
\item **Usos:** converter entre densidade e pressão/temperatura reais; calcular o fator de volume $B_g$; avaliar compressibilidade isotérmica do gás; PVT e balanço de massas.
\end{itemize}

---

\textbf{1) Formulação e relações básicas}

\begin{itemize}
\item Molar volume: $V_m = \dfrac{Z R T}{p}$.  
\item Densidade (massa): $\rho = \dfrac{M p}{Z R T}$, onde $M$ é massa molar (kg·mol^{-1}).
\item Fator de volume do gás (molar): $B_{m} = V_m = \dfrac{Z R T}{p}$. (Em reservatórios usa‑se versões por massa/por scf — adaptar unidades.)
\end{itemize}

\textbf{2) Relação com compressibilidade isotérmica do gás $c_g$}

Partindo de $V_m = Z R T / p$ e definindo $c_g = -\dfrac{1}{V_m} \left( \dfrac{\partial V_m}{\partial p} \right)_T$ obtém‑se:

$$
c_g \,=\, \frac{1}{p} \, - \, \frac{1}{Z} \left(\frac{\partial Z}{\partial p}\right)_T \,=\, \frac{1}{p} \, - \, \frac{\partial \ln Z}{\partial p} .
$$

Interpretação prática: para gás quase ideal $\partial Z/\partial p \approx 0$ e $c_g\approx 1/p$. O termo com $\partial Z/\partial p$ corrige o comportamento real; portanto é importante para pressões altas / não‑ideais.

\textbf{3) Métodos para obter $Z$ (ordem de complexidade e precisão)}

\begin{itemize}
\item Leitura direta no gráfico Standing–Katz (rápido; precisa dos pseudo‑críticos do gás).  
\item Correlações empíricas e semi‑empíricas (ex.: Dranchuk–Abou‑Kassem) — boas para implementação numérica.  
\item Equações de estado cúbicas (Peng–Robinson, Soave–Redlich–Kwong) — recomendadas quando se tem composição e se precisa de consistência em fase líquida/vapor; permitem calcular Z e derivadas termodinâmicas.
\end{itemize}

\textbf{4) Pseudocríticos e preparação (how‑to)}

Para usar chart ou correlações reduzidas, calcule $T_{pc}$ e $P_{pc}$ da mistura. Duas rotas:

\begin{itemize}
\item Regra de Kay (mistura):  
\item para cada componente i obtenha $T_{c,i}$ e $P_{c,i}$ (temperaturas Kelvin, pressão em Pa).  
\item calcule $T_{pc} = \sum_i y_i \, T_{c,i}$ e $P_{pc} = \sum_i y_i \, P_{c,i}$ (mole fractions $y_i$).  
\item então $T_r = T / T_{pc}$ e $P_r = p / P_{pc}$.  
\end{itemize}

\begin{itemize}
\item Correções para gás ácido / CO2/H2S (Wichert–Aziz): existem procedimentos para reduzir $T_{pc}$ e $P_{pc}$ quando frações significativas de CO2/H2S estiverem presentes — aplicar se $\mathrm{CO_2}+\mathrm{H_2S}$ > ~2–5\% mole.
\end{itemize}

Com $T_r$ e $P_r$ você lê $Z$ no gráfico de Standing–Katz (curvas $Z$ vs $P_r$ para cada $T_r$). Para automação prefira correlações numéricas (Dranchuk–Abou‑Kassem) ou uma EoS.

\textbf{5) Método robusto: Peng–Robinson (PR) — passo a passo (Padrão industrial)}

Parâmetros do PR (unidades SI):

$$
a = 0.45724 \frac{R^2 T_c^2}{P_c},\qquad b = 0.07780 \frac{R T_c}{P_c}
$$

onde $R=8.314462618$ J·mol^{-1}·K^{-1}, $T_c$ (K), $P_c$ (Pa). O fator $\alpha$ descreve dependência com $T$:

$$
\kappa = 0.37464 + 1.54226\,\omega - 0.26992\,\omega^2,\qquad
\alpha(T) = \left(1 + \kappa \left(1-\sqrt{\dfrac{T}{T_c}}\right)\right)^2
$$

Defina então:

$$
A = \dfrac{a\,\alpha(T)\,p}{R^2 T^2},\qquad B = \dfrac{b p}{R T}.
$$

O polinômio cúbico em $Z$ (forma padrão PR):

$$
Z^3 - (1 - B) Z^2 + (A - 3B^2 - 2B) Z - (A B - B^2 - B^3) = 0.
$$

Solução: encontre as raízes reais; para fase gasosa escolha a maior raiz real (maior $Z$).

\textbf{6) Exemplo numérico completo (CH4 puro)}

Dados (SI):
\begin{itemize}
\item $T = 350\ \mathrm{K}$ (exemplo)
\item $p = 5\ \mathrm{MPa} = 5\times10^6\ \mathrm{Pa}$
\item Metano: $T_c = 190.564\ \mathrm{K}$, $P_c = 4.5992\times10^6\ \mathrm{Pa}$, $\omega\approx0.011$.
\end{itemize}

\begin{enumerate}
\item Calcule $a,b$:
\end{enumerate}

$$
a \approx 0.45724\dfrac{R^2 T_c^2}{P_c} \approx 0.2498\quad(\mathrm{SI})
$$
$$
b \approx 0.07780\dfrac{R T_c}{P_c} \approx 2.68\times10^{-5}\ \mathrm{m^3\,mol^{-1}}
$$

\begin{enumerate}
\item $\kappa\approx 0.3916$, $\alpha\approx 0.7396$ (usar fórmula acima).
\end{enumerate}

\begin{enumerate}
\item $A\approx 0.1089$, $B\approx 0.0460$ (ver cálculos no script abaixo).
\end{enumerate}

\begin{enumerate}
\item Monta o cúbico e resolve: obtém $Z\approx 0.948$ (raiz gasosa principal).
\end{enumerate}

Interpretação: $Z<1$ porém próximo de 1 — desvio moderado para estas condições.

\textbf{7) Compressibilidade isotérmica calculada numericamente (exemplo)}

Procedimento prático (numérico):
\begin{enumerate}
\item Calcule $Z(p)$ com o método EoS (PR) no ponto $p$ (aqui 5 MPa).
\item Calcule $Z(p+\Delta p)$ com $\,\Delta p$ pequeno (ex.: $10^4$–$10^5\ \mathrm{Pa}$).  
\item Estime $\partial Z/\partial p \approx (Z(p+\Delta p)-Z(p))/\Delta p$.  
\item Use $c_g = 1/p - (1/Z) (\partial Z/\partial p)$.
\end{enumerate}

Usando o exemplo numérico (valores aproximados do passo anterior e $\Delta p=10^5\ \mathrm{Pa}$) obtemos $c_g\approx 2.19\times10^{-7}\ \mathrm{Pa^{-1}}$ que corresponde a $0.219\ \mathrm{MPa^{-1}}$ ou aproximadamente $1.51\times10^{-3}\ \mathrm{psi^{-1}}$ (conversões mostradas no script).

\textbf{8) Código Python (PR EoS) — cálculo de Z e compressibilidade por diferença finita}

Instalar dependência mínima: \texttt{pip install numpy}

\begin{verbatim}
import numpy as np

R = 8.314462618

def Z_PR(p, T, Tc, Pc, omega):
    # p in Pa, T in K, Tc in K, Pc in Pa
    a = 0.45724 * R**2 * Tc**2 / Pc
    b = 0.07780 * R * Tc / Pc
    kappa = 0.37464 + 1.54226*omega - 0.26992*omega**2
    alpha = (1 + kappa*(1 - np.sqrt(T/Tc)))**2
    A = a * alpha * p / (R**2 * T**2)
    B = b * p / (R * T)
    # Coeffs cubic: Z^3 - (1-B) Z^2 + (A - 3B^2 - 2B) Z - (A B - B^2 - B^3) = 0
    coeffs = [1.0, -(1 - B), (A - 3*B*B - 2*B), -(A*B - B*B - B*B*B)]
    roots = np.roots(coeffs)
    real_roots = np.real(roots[np.isclose(roots.imag, 0, atol=1e-8)])
    if len(real_roots) == 0:
        raise RuntimeError('No real root found for Z')
    # gas root = largest real root
    Z = np.max(real_roots)
    return float(Z)

def gas_compressibility(p, T, Tc, Pc, omega, dp=1e5):
    Z1 = Z_PR(p, T, Tc, Pc, omega)
    Z2 = Z_PR(p + dp, T, Tc, Pc, omega)
    dZdp = (Z2 - Z1) / dp
    c_g = 1.0/p - (1.0/Z1)*dZdp
    return Z1, dZdp, c_g

if __name__ == '__main__':
    # Example: methane at p=5 MPa, T=350 K
    Tc = 190.564
    Pc = 4.5992e6
    omega = 0.011
    p = 5e6
    T = 350.0
    Z, dZdp, cg = gas_compressibility(p, T, Tc, Pc, omega, dp=1e5)
    print('Z =', Z)
    print('dZ/dp =', dZdp, 'Pa^-1')
    print('c_g =', cg, 'Pa^-1 (', cg*1e6, 'MPa^-1 )')
\end{verbatim}

Saída esperada aproximada (exemplo que usamos manualmente): \texttt{Z ≈ 0.948}, \texttt{c_g ≈ 2.19e-7 Pa^-1}.

\textbf{9) Recomendações práticas e atenção a unidades}

\begin{itemize}
\item Sempre usar unidades consistentes (SI: Pa, K, J/(mol·K)). Para conversões (psi, °R) ajuste constantes do gás.  
\item Para misturas use regras de mistura apropriadas (EoS: combinações de `a_{ij}`/fatores de interação). Se não conhecer os coeficientes binários use $k_{ij}=0$ como primeira aproximação (pode introduzir erro).  
\item Para trabalho de campo rápido: calculadoras com Standing–Katz (chart) ou correlação DAK. Para simulação/reconhecimento e prognóstico, prefira EoS com mistura.
\end{itemize}

\textbf{10) Referências rápidas}

\begin{itemize}
\item Standing, M. B., Katz, D. L. — gráfico Standing–Katz para $Z$ (uso clássico).  
\item Dranchuk, P. M. & Abou‑Kassem, J. H. (1975) — correlação numérica para $Z$ (boa para automação).  
\item Peng, D.-Y. & Robinson, D. B. (1976) — Peng–Robinson EoS (cúbica), uso industrial.  
\item Wichert, D. & Aziz, K. (1972) — correção pseudo‑crítica para CO2/H2S (gases ácidos).
\end{itemize}

---

Se quiser, eu:
\begin{itemize}
\item gero uma versão com figuras (Standing–Katz) e tabela de propriedades críticas;  
\item adiciono uma implementação DAK (Dranchuk–Abou‑Kassem) completa em Python;  
\item aplico o cálculo a uma mistura real (envie composição molar) e salvo o resultado.
\end{itemize}

% ---- file: exercicios_resolvidos.md
\section{Exercícios Resolvidos — Engenharia de Reservatórios I}

Compilação de soluções passo‑a‑passo extraídas dos ficheiros em \texttt{Estudar/matéria} (cap1–cap5). Use como referência rápida para estudo; ver ficheiros de origem para enunciados completos.

---

\subsection{Capítulo 1 — Exemplo rápido (OOIP)}
Dados: Área = 100 acres; $h=20\,$ft; $\phi=0.15$; $S_{wi}=0.25$; $B_o=1.2\,$bbl/STB.

Fórmula prática:
$$OOIP=\dfrac{7758\,A\,h\,\phi\,(1-S_{wi})}{B_o}$$

Cálculo passo a passo:
\begin{align*}
7758\times100&=775\,800\\
775\,800\times20&=15\,516\,000\\
15\,516\,000\times0.15&=2\,327\,400\\
2\,327\,400\times0.75&=1\,745\,550\\
OOIP&=\dfrac{1\,745\,550}{1.2}\approx1\,454\,625\;\text{STB}
\end{align*}

Resultado: ≈1.455×10^6 STB.

---

\subsection{Capítulo 2 — PVT: cálculo de $B_o$ e conversão de $R_s$}
Enunciado (exemplo): 1 STB amostrado dá $V_{res}=1.20\,$bbl, $V_{surf}=1.00\,$STB; $R_s=400\,$scf/STB.

\begin{enumerate}
\item $B_o$:
$$B_o=\dfrac{V_{res}}{V_{surf}}=\dfrac{1.20}{1.00}=1.20\;\text{bbl/STB}.$$ 
\end{enumerate}

\begin{enumerate}
\item Conversão de $R_s$ para m^3/m^3 (opcional):
1 scf = 0.0283168 m^3; 1 STB = 0.1589873 m^3.
$$R_s(\mathrm{m}^3/\mathrm{m}^3)=\dfrac{400\times0.0283168}{0.1589873}\approx71.3\;\mathrm{m}^3/\mathrm{m}^3.$$ 
\end{enumerate}

Interpretação: $B_o>1$ indica expansão/volume maior em reservatório por unidade de superfície; $R_s$ alto = muito gás dissolvido.

---

\subsection{Capítulo 3 — Rochas: porosidade (massa e Arquímedes)}

\subsubsection{Exemplo 1 (Exercício 1)}
Dados: $m_{sat}=130\,$g; $m_{dry}=105\,$g; $\rho_o=0.84\,$g/cm^3; $V_t=180\,$cm^3.

\begin{enumerate}
\item Volume de fluido nos poros:
$$V_f=\dfrac{m_{sat}-m_{dry}}{\rho_o}=\dfrac{130-105}{0.84}=\dfrac{25}{0.84}\approx29.7619\;\text{cm}^3.$$ 
\end{enumerate}

\begin{enumerate}
\item Porosidade:
$$\phi=\dfrac{V_p}{V_t}=\dfrac{29.7619}{180}\approx0.16534\approx16.53\%.$$ 
\end{enumerate}

\subsubsection{Exemplo 2 (método de Arquímedes — Exercício 5)}
Dados: $m_{dry}=330\,$g; $m_{sat}=360\,$g; $m_{ap\_agua}=225\,$g; $\rho_{agua}=1\,$g/cm^3.

\begin{enumerate}
\item Volume total via empuxo:
$$V_t=\dfrac{m_{sat}-m_{ap\_agua}}{\rho_{agua}}=\dfrac{360-225}{1}=135\;\text{cm}^3.$$ 
\end{enumerate}

\begin{enumerate}
\item Volume poroso:
$$V_p=m_{sat}-m_{dry}=360-330=30\;\text{cm}^3.$$ 
\end{enumerate}

\begin{enumerate}
\item Porosidade:
$$\phi=\dfrac{V_p}{V_t}=\dfrac{30}{135}\approx0.22222\approx22.22\%.$$
\end{enumerate}

Observação: indicar unidades e arredondamentos no relatório final.

---

\subsection{Capítulo 4 — Cálculo volumétrico (Exercício 4.8 — OOIP)}
Dados: $A=200\,$acres; $h_{net}=30\,$ft; $\phi=0.18$; $S_{wi}=0.25$; $B_o=1.2\,$bbl/STB.

Fórmula:
$$OOIP=\dfrac{7758\,A\,h\,\phi\,(1-S_w)}{B_o}.$$ 

Cálculo:
\begin{align*}
7758\times200&=1\,551\,600\\
1\,551\,600\times30&=46\,548\,000\\
46\,548\,000\times0.18&=8\,378\,640\\
8\,378\,640\times0.75&=6\,283\,980\\
OOIP&=\dfrac{6\,283\,980}{1.2}\approx5\,236\,650\;\text{STB}.
\end{align*}

Resultado: ≈5.24×10^6 STB.

---

\subsection{Capítulo 5 — Equação de Balanço de Materiais (Havlena & Odeh)}

\subsubsection{Exemplo (cálculo do termo $E_o$)}
Dados ilustrativos: $B_{o,i}=1.10$, $B_o=1.15$, $R_{s,i}=200\,$scf/STB, $R_s=180\,$scf/STB, $B_g=0.005$.

Fórmula:
$$E_o = B_o - B_{o,i} + (R_{s,i}-R_s)B_g$$

Cálculo:
$$E_o = 1.15 - 1.10 + (200-180)\times0.005 = 0.05 + 20\times0.005 = 0.05 + 0.10 = 0.15.$$ 

Uso: repita para cada data (pressão) e monte a regressão linear $F= N E_o + m E_g + (1+m)E_{f,w} + W_e$ para estimar $N$ e $m$.

---

\subsection{Observações finais}
\begin{itemize}
\item Este ficheiro agrupa exemplos resolvidos presentes nas notas (cap1–cap5). Para exercícios completos e figuras, consulte os ficheiros originais em `Estudar/matéria`.
\item Se quiser, posso:
\item adicionar soluções passo‑a‑passo para TODOS os exercícios listados nos capítulos (criando uma secção por número de exercício), ou
\item gerar um PDF/LaTeX com este compêndio.
\end{itemize}

---

\emph{Ficheiro gerado automaticamente pelo assistente.}

---

\subsection{Soluções adicionais completas — Capítulos 1–5}

A seguir estão as soluções passo‑a‑passo para as perguntas listadas nas secções de consolidação e tarefas dos Capítulos 1 a 5. Para problemas descritivos forneço uma resposta‑modelo; para exercícios numéricos apresento cálculos explícitos e resultados.

\subsubsection{Capítulo 1 — Respostas e justificações (Consolidação)}
\begin{enumerate}
\item **B** — Baixa densidade e alta mobilidade.
\begin{itemize}
\item Justificação: óleo leve tem menor densidade e viscosidade, o que aumenta a mobilidade relativa e facilita o escoamento.
\end{itemize}
\end{enumerate}

\begin{enumerate}
\item **A** — Reduz a pressão do reservatório e facilita o fluxo de óleo.
\begin{itemize}
\item Justificação: a presença de gás (dissolvido ou livre) altera o comportamento volumétrico e pode gerar mecanismos de drive (solution gas, gas‑cap) que afetam a pressão e a mobilidade do óleo.
\end{itemize}
\end{enumerate}

\begin{enumerate}
\item **A** — Manutenção da pressão por influxo de água do aquífero.
\begin{itemize}
\item Justificação: water drive é caracterizado pelo suporte de pressão fornecido pelo aquífero que desloca óleo para os poços.
\end{itemize}
\end{enumerate}

\begin{enumerate}
\item **A** — A taxa de fluxo pelo reservatório.
\begin{itemize}
\item Justificação: a mobilidade é proporcional a $1/\mu$; aumentos de viscosidade reduzem caudal para a mesma diferença de pressão.
\end{itemize}
\end{enumerate}

\begin{enumerate}
\item **A** — Mantém pressão por expansão de gás acima do óleo.
\begin{itemize}
\item Justificação: gas‑cap drive usa a expansão do gás livre para suportar pressão de reservatório.
\end{itemize}
\end{enumerate}

\begin{enumerate}
\item **A** — O condensado separa‑se do gás à medida que a pressão diminui.
\begin{itemize}
\item Justificação: condensado aparece quando a condição de orvalho é atingida e parte do gás condensa em fase líquida.
\end{itemize}
\end{enumerate}

\begin{enumerate}
\item **E** — (não define).  
\begin{itemize}
\item Explicação: as características A–D são diretamente associadas ao solution‑gas drive; a expressão "produção contínua de óleo" não é uma definição intrínseca do mecanismo (a produção pode declinar com a queda de pressão), portanto é a alternativa que NÃO define o mecanismo.
\end{itemize}
\end{enumerate}

\begin{enumerate}
\item **A** — A água do aquífero desloca o óleo em direção aos poços.
\end{enumerate}

\begin{enumerate}
\item **A** — Pressão inicial do reservatório e mobilidade do óleo.
\begin{itemize}
\item Justificação: em reservatórios undersaturated a produção depende da pressão acima do ponto de bolha e da mobilidade do óleo.
\end{itemize}
\end{enumerate}

\begin{enumerate}
\item **A** — Gás livre no topo do reservatório que ajuda a manter a pressão.
\end{enumerate}

---

\subsubsection{Capítulo 2 — Propriedades dos fluidos (respostas‑modelo)}
\begin{itemize}
\item O que é $B_o$ e por que é importante:
\item $B_o$ é o fator de volume de formação do óleo: $B_o=V_{res}/V_{surf}$. É usado para converter volumes no reservatório para volumes de superfície (STB) e é essencial em balanços de massa e estimativas de reservas.
\end{itemize}

\begin{itemize}
\item Como $R_s$ varia com pressão e significado do ponto de bolha:
\item Geralmente $R_s$ diminui com a redução de pressão; ao atingir o ponto de bolha surge gás livre e $R_s$ passa a ser o valor máximo de gás dissolvido na condição dada.
\end{itemize}

\begin{itemize}
\item Efeito da viscosidade do óleo na mobilidade:
\item Mobilidade dinâmica $\\lambda = k/\\mu$; viscosidade maior reduz mobilidade e diminui vazões para um mesmo gradiente de pressão.
\end{itemize}

\begin{itemize}
\item Quando tratar o gás como ideal vs real:
\item Use o modelo ideal quando as condições estiverem longe do crítico (pressões relativamente baixas e temperaturas altas); para pressões elevadas/temperaturas baixas use fator $Z$ (gráficos de Standing & Katz) ou equações de estado cúbicas (Peng‑Robinson, SRK).
\end{itemize}

---

\subsubsection{Capítulo 3 — Rochas (soluções numéricas e procedimentos)}

\begin{enumerate}
\item (já incluído) — Porosidade por pesagem: exemplo resolvido no ficheiro acima (Exemplo 1).
\end{enumerate}

\begin{enumerate}
\item (figuras) — Porosidade idealizada: desenhe o volume total $V_t$, identifique $V_p$ (vazios) e calcule $\phi=V_p/V_t$. Em figuras, calcule áreas/volumes preferenciais e aplique fórmula.
\end{enumerate}

\begin{enumerate}
\item **Exercício (válvula e pressões)** — Dados: $V_1=100\,$cc; $V_2=100\,$cc; $p_1=15\,$psi; $p_2=60\,$psi; $p_f=39\,$psi. Calcular volume do grão $V_g$.
\end{enumerate}

Derivação (conservação do número de moles a temperatura constante — Lei de Boyle):
$$p_1(V_1-V_g)+p_2V_2=p_f\,(V_1-V_g+V_2).$$
Resolvendo para $V_g$:
$$V_g=\dfrac{p_f(V_1+V_2)-p_1V_1-p_2V_2}{p_f-p_1}.$$ 
Substituindo valores:
$$V_g=\dfrac{39(100+100)-15\times100-60\times100}{39-15}=\dfrac{7800-1500-6000}{24}=\dfrac{300}{24}=12{,}5\;\text{cm}^3.$$

Resposta: $V_g=12{,}5\,$cc.

\begin{enumerate}
\item **Porosidade média (Exercício 4)** — dados: 10, 12, 11, 13, 14, 10, 17 (%).
\begin{itemize}
\item Soma = 77; média = $77/7 = 11.0\%$.
\end{itemize}
\end{enumerate}

\begin{enumerate}
\item (já incluído) — Método de Arquímedes: exemplo resolvido no ficheiro (Exemplo 2).
\end{enumerate}

\begin{enumerate}
\item **Compressibilidade de poros (exemplo)** — Dados: $V_p=18\,$cm^3$; \Delta V_p=0{,}15\,$cm^3 para $\Delta p=900\,$psi.
$$C_f=\dfrac{\Delta V_p/V_p}{\Delta p}=\dfrac{0{,}15/18}{900}\approx9{,}26\times10^{-6}\;\text{psi}^{-1}.$$ 
\end{enumerate}

---

\subsubsection{Capítulo 4 — Cálculo volumétrico (sensibilidade e procedimentos)}

**Exercício 4.8 (já incluído)** — OOIP calculado: aproximadamente $5{,}236{,}650\,$STB para $A=200\,$acres, $h=30\,$ft, $\phi=0.18$, $S_{wi}=0.25$, $B_o=1.2$.

**Análise de sensibilidade (variação de $\phi$ de 0{,}15 a 0{,}22)** — usando a mesma fórmula prática:
\begin{itemize}
\item Para $\phi=0{,}15$:
$$OOIP\approx4{,}363{,}875\;\text{STB}.$$ 
\item Para $\phi=0{,}22$:
$$OOIP\approx6{,}400{,}350\;\text{STB}.$$ 
\end{itemize}

Interpretação: OOIP cresce aproximadamente linearmente com $\phi$ (mesmo A,h,Swi e Bo), mostrando elevada sensibilidade à porosidade; variação relativa de $\phi$ traduz‑se diretamente numa variação proporcional do OOIP.

\textbf{Procedimento Monte Carlo (resumo prático):}
\begin{itemize}
\item Escolher distribuições (ex.: normal/triangular) para $\phi$, $S_w$, $B_o$, $h$.
\item Gerar N iterações (p.ex. 1.000–10.000), calcular OOIP por iteração e recolher estatísticas P10/P50/P90.
\item Reportar média, desvios e percentis; preparar gráficos de dispersão e histogramas.
\end{itemize}

---

\subsubsection{Capítulo 5 — Equação de Balanço de Materiais (EBM) — passos e exemplo}

\textbf{Passos para aplicar a linearização (Havlena \& Odeh):}
\begin{enumerate}
\item Reunir por data: $p$, $N_p$, $G_p$, $W_p$ (produções), $B_o(p)$, $B_g(p)$, $R_s(p)$ a partir de PVT.
\item Calcular colunas por data: $E_o,E_g,E_{f,w}$ usando as fórmulas:
$$E_o = B_o - B_{o,i} + (R_{s,i}-R_s)B_g$$
$$E_g = B_{o,i}\left(\dfrac{B_g}{B_{g,i}} - 1\right)$$
$$E_{f,w} = B_{o,i}\left(S_{w,i}c_w + \dfrac{c_f}{1-S_{w,i}}\Delta p\right)$$
\item Montar $F$ (lado conhecido) por data (converter unidades):
$$F = N_p B_o + G_p - R_s B_g + W_p B_w - W_{inj} B_w - G_{inj} B_{g,inj}.$$ 
\item Fazer regressão linear múltipla: $F = N E_o + m E_g + (1+m) E_{f,w} + W_e$ para estimar $N$ (estoque original) e $m$ (razão gas‑cap/óleo). Use regressão com todas as linhas de dados (mínimo 3–4 pontos; quanto mais, melhor).
\end{enumerate}

\textbf{Exemplo (termo $E_o$ já calculado no ficheiro):}
\begin{itemize}
\item Valores ilustrativos: $B_{o,i}=1{,}10$, $B_o=1{,}15$, $R_{s,i}=200\,$scf/STB, $R_s=180\,$scf/STB, $B_g=0{,}005$.
\item Calculámos:
$$E_o=1{,}15-1{,}10+(200-180)\times0{,}005=0{,}15.$$ 
\end{itemize}

Para obter $N$ proceda assim (resumo): construa a tabela com colunas $F,E_o,E_g,E_{f,w}$ e execute uma regressão linear múltipla (p.ex. em Excel: Análise de Regressão → regressão múltipla; em Python: numpy.linalg.lstsq ou statsmodels.OLS). O coeficiente associado a $E_o$ será a estimativa de $N$ (ou, dependendo da forma, a sua relação direta — ver convenção de montagem das colunas).

---

Se desejar, posso agora:
\begin{itemize}
\item (A) inserir soluções detalhadas (passo‑a‑passo) para cada exercício numerado nos ficheiros originais (por exemplo, 1–N de `exercícios_transcription.md`), incluindo os cálculos intermédios e tabelas; ou
\item (B) gerar um PDF/LaTeX do compêndio completo `exercicios_resolvidos.md` com formatação académica.
\end{itemize}

Ficarei aguardando a sua preferência para a etapa seguinte.

% ---- file: exame_gabarito.md
\section{Revisão para a prova — Gabarito e Rubricas}

Este ficheiro reúne as 4 questões propostas para a prova, as respostas‑gabarito e as rubricas de correção.

---

\subsection{Questão 1 — Verdadeiro ou Falso (Capítulo 1)}

Enunciado (10 itens, 1 pt cada):
\begin{enumerate}
\item A armadilha geológica impede a migração de hidrocarbonetos e permite sua acumulação.
\item Rocha geradora é caracterizada por baixa matéria orgânica.
\item Rocha selante tem alta permeabilidade e facilita o fluxo de fluidos.
\item Migração de hidrocarbonetos ocorre preferencialmente por meios porosos e fraturados.
\item Sincronismo (coincidência temporal de processos) é irrelevante para a formação de um sistema petrolífero.
\item Catagênese refere‑se ao craqueamento térmico do querogênio em hidrocarbonetos.
\item O sistema de produção inclui instalações de coleta, elevação e separação até a superfície.
\item Em um gas‑cap drive, o gás livre contribui para manter a pressão do reservatório.
\item Porosidade e permeabilidade são a mesma propriedade física.
\item O fator de volume de formação `B_o` relaciona volumes em reservatório ao volume na superfície.
\end{enumerate}

\subsubsection{Gabarito Q1}
\begin{enumerate}
\item V
\item F
\item F
\item V
\item F
\item V
\item V
\item V
\item F
\item V
\end{enumerate}

\subsubsection{Rubrica Q1 (10 pts)}
\begin{itemize}
\item 1 ponto por resposta correta.
\item Se for solicitada justificativa: aceitar 1–2 frases; atribuir 0.5 pt para justificativa parcial.
\end{itemize}

---

\subsection{Questão 2 — Dissertação (Capítulo 4)}

Enunciado:
Redija uma dissertação (≈ 400–600 palavras) que responda e discuta:
\begin{itemize}
\item Descreva o método volumétrico para estimativa de OOIP/OGIP e os principais parâmetros envolvidos (`A`, `h`, `\phi`, `S_w`, `B_o`, `B_g`).
\item Explique as principais fontes de incerteza (cut‑offs, N/G, heterogeneidade, PVT) e como afetam a estimativa.
\item Proponha um fluxo de trabalho prático para integrar mapas de net, dados de core, logs e PVT para obter uma estimativa com P10/P50/P90 (mencione Monte Carlo).
\item Conclua com recomendações para comunicar risco e medidas para reduzir incerteza (QA/QC, aquisição de dados adicionais, sensibilidade).
\end{itemize}

\subsubsection{Rubrica Q2 (30 pts)}
\begin{itemize}
\item Método volumétrico e parâmetros (12 pts): fórmula prática (ex.: $$OOIP=\dfrac{7758\,A\,h\,\phi\,(1-S_w)}{B_o}$$), explicação de cada termo, menção a unidades campo vs SI. (12 pts)
\item Fontes de incerteza (8 pts): identificação e explicação do impacto de cut‑offs, N/G, heterogeneidade, erros PVT. (8 pts)
\item Fluxo de trabalho prático (6 pts): integração cores/logs/PVT, construção de mapas, cálculo setorial, sensibilidade e Monte Carlo para P10/P50/P90. (6 pts)
\item Comunicação de risco e recomendações (4 pts): QA/QC, aquisição adicional, apresentação de Pxx e mitigação. (4 pts)
\end{itemize}

Critérios de avaliação:
\begin{itemize}
\item Conteúdo técnico: 60% dos pontos.
\item Metodologia/prática: 20%.
\item Comunicação/clareza e recomendações: 20%.
\end{itemize}

---

\subsection{Questão 3 — Cálculo / Prático (Capítulo 2 — PVT)}

Enunciado:
Um ensaio PVT fornece os seguintes resultados para 1 STB de óleo amostrado: volume no estado de reservatório $V_{res}=1.20\,$bbl; volume correspondente à superfície $V_{surf}=1.00\,$STB; gás libertado no ensaio $R_s=400\,$scf/STB.

\begin{enumerate}[label=\Alph*)]
\item Calcule o fator de volume de formação `B_o` (bbl/STB) usando $B_o=V_{res}/V_{surf}$.\n\end{enumerate}

\begin{enumerate}[label=\Alph*)]
\item Expresse `R_s` em m³/m³ usando $1\,\text{scf}=0.0283168\,\text{m}^3$ e $1\,\text{STB}=0.1589873\,\text{m}^3$.\n\end{enumerate}

\begin{enumerate}[label=\Alph*)]
\item Interprete fisicamente um $B_o>1$ e discuta brevemente como $R_s$ influencia o comportamento de produção.\n\end{enumerate}

\subsubsection{Gabarito Q3 (30 pts — 10/10/10)}
(a) $$B_o=\dfrac{V_{res}}{V_{surf}}=\dfrac{1.20}{1.00}=1.20\;\text{bbl/STB}.$$ (10 pts)

(b) Conversão: $$R_s\,\text{(m}^3/\text{m}^3)=\dfrac{400\times0.0283168}{0.1589873}\approx71.3\;\text{m}^3/\text{m}^3.$$ (10 pts)

(c) Interpretação: (10 pts)
\begin{itemize}
\item $B_o>1$ indica que 1 STB à superfície corresponde a mais de 1 bbl nas condições de reservatório (efeito de compressibilidade e/ou gás dissolvido); volumes no reservatório são maiores que o volume final surfacel.
\item $R_s$ alto implica presença significativa de gás dissolvido; durante declínio de pressão, se atingir o ponto de bolha, o gás é liberado, alterando mobilidade do fluido, afetando recuperação e mecanismo de produção. Pontos por clareza e ligação a efeitos operacionais.
\end{itemize}

---

\subsection{Questão 4 — Cálculo / Prático (Capítulo 3 — Rocha)}

Enunciado:
\begin{enumerate}[label=\Alph*)]
\item Porosidade por método de massa: uma amostra tem massa seca $m_{dry}=250\,$g e massa saturada $m_{sat}=290\,$g; densidade do fluido $\rho_f=0.90\,$g/cm³; volume total da amostra $V_t=320\,$cm³. Calcule a porosidade $\phi$ (use $V_p=(m_{sat}-m_{dry})/\rho_f$ e $\phi=V_p/V_t$).\n\end{enumerate}

\begin{enumerate}[label=\Alph*)]
\item Permeabilidade (Darcy): num ensaio de bancada obteve‑se vazão volumétrica $q=1.0\times10^{-6}\,\text{m}^3/s$, área da amostra $A=10\,\text{cm}^2=1.0\times10^{-3}\,\text{m}^2$, comprimento $L=5\,\text{cm}=0.05\,\text{m}$, viscosidade $\mu=1.0\,\text{cP}=0.001\,\text{Pa·s}$ e queda de pressão $\Delta p=10{,}000\,$Pa. Calcule `k` em m² e converta para Darcy (1 D = 9.869233×10⁻¹³ m²) usando a relação:\n$$k=\dfrac{q\,\mu\,L}{A\,\Delta p}$$
\end{enumerate}

\subsubsection{Gabarito Q4 (30 pts — 15/15)}
(a) Porosidade (15 pts):
\begin{itemize}
\item $$V_p=\dfrac{m_{sat}-m_{dry}}{\rho_f}=\dfrac{290-250}{0.90}=\dfrac{40}{0.90}=44.444\;\text{cm}^3.$$ 
\item $$\phi=\dfrac{V_p}{V_t}=\dfrac{44.444}{320}\approx0.1389\approx13.9\%. $$
(atribuir parcial se passos corretos; tolerância ±0.2–0.5\% absoluta)
\end{itemize}

(b) Permeabilidade (15 pts):
\begin{itemize}
\item Aplicando Darcy:
  $$k=\dfrac{1\times10^{-6}\times0.001\times0.05}{1\times10^{-3}\times10{,}000}=5\times10^{-12}\;\text{m}^2.$$ 
\item Converter para Darcy:
  $$k_{D}=\dfrac{5\times10^{-12}}{9.869233\times10^{-13}}\approx5.07\;\text{D}=5\,070\;\text{mD}.$$ 
(Aceitar pequenas variações por arredondamento; exigir todas as conversões)
\end{itemize}

---

\subsection{Pontuação total sugerida}
\begin{itemize}
\item Q1: 10 pts
\item Q2: 30 pts
\item Q3: 30 pts
\item Q4: 30 pts
\end{itemize}

Total: 100 pts.

\subsection{Observações finais sobre correção}
\begin{itemize}
\item Para cálculos, exigir passos e unidades; penalizar falta de conversões ou unidades incoerentes.
\item Para dissertação, pontuar conteúdo técnico, metodologia e clareza.
\item Para todos os itens, aceitar pequenas variações numéricas quando justificadas por arredondamento.
\end{itemize}

---

\emph{Arquivo gerado automaticamente pelo assistente.}
\end{document}\n